\documentclass{ieeetj}
\usepackage{cite}
\usepackage{amsmath,amssymb,amsfonts}
\usepackage{algorithmic}
\usepackage{graphicx,color}
\usepackage{textcomp}
\usepackage{xcolor}
\usepackage{hyperref}
\usepackage{xcolor}
\usepackage{longtable}
\usepackage{booktabs}
\usepackage{array}
\usepackage{caption}
\usepackage{ragged2e}
\usepackage{svg}
\usepackage{enumitem}
\hypersetup{hidelinks=true}
\usepackage{algorithm,algorithmic}
\usepackage{multirow}
\def\BibTeX{{\rm B\kern-.05em{\sc i\kern-.025em b}\kern-.08em
    T\kern-.1667em\lower.7ex\hbox{E}\kern-.125emX}}
\AtBeginDocument{\definecolor{tmlcncolor}{cmyk}{0.93,0.59,0.15,0.02}\definecolor{NavyBlue}{RGB}{0,86,125}}

\newcommand{\mfabric}{M\textsuperscript{2}Fabric}

\def\OJlogo{\vspace{-4pt}$<$Society logo(s) and publication title will appear here.$>$}
\def\seclogo{\vspace{10pt}$<$Society logo(s) and publication title will appear here.$>$}

\def\authorrefmark#1{\ensuremath{^{\textbf{#1}}}}

\begin{document}
\receiveddate{XX Month, XXXX}
\reviseddate{XX Month, XXXX}
\accepteddate{XX Month, XXXX}
\publisheddate{XX Month, XXXX}
\currentdate{XX Month, XXXX}
\doiinfo{XXXX.2022.1234567}

\markboth{}{Author {et al.}}

\title{\mfabric{}: A Multilayer Post-Quantum Cryptographic Architecture for Hyperledger Fabric}

\author{Zara Laila Cheong\authorrefmark{1}, Ali Tufail\authorrefmark{1}, and Rosyzie Anna Apong\authorrefmark{1}}
\affil{School of Digital Science, Universiti Brunei Darussalam, Bandar Seri Begawan, BE1410, Brunei Darussalam}
\corresp{Corresponding author: Zara Laila Cheong (email: zaralaila@icloud.com).}
\authornote{The authors declare no specific funding for this work and no conflicts of interest. The implementation, benchmark scripts, and raw measurement data are openly available in the M\textsuperscript{2}Fabric repository to support independent reproduction.}


\begin{abstract}
{Hyperledger Fabric relies on classical elliptic-curve cryptography across its
entire transaction lifecycle and identity infrastructure, making every signing,
certification, validation, and key-exchange operation vulnerable to a
cryptographically relevant quantum computer. Prior post-quantum proposals for
Fabric have each addressed only a subset of its cryptographic layers, leaving
the unaddressed layers quantum-exploitable. This article presents \mfabric{}, a
multilayer post-quantum architecture that replaces all four live cryptographic
layers of Hyperledger Fabric v2.5.12 and Fabric CA v1.5.12 simultaneously: the
Blockchain Cryptographic Service Provider (BCCSP) signing layer, the
certificate infrastructure, the Membership Service Provider
identity-validation layer, and the Transport Layer Security (TLS) inter-node
transport layer. The design uses the NIST-standardised algorithms ML-DSA-44
(FIPS~204, Category~2), ML-DSA-65 (FIPS~204, Category~3), FN-DSA-512
(FIPS~206 Initial Public Draft, Category~1), and ML-KEM-768 (FIPS~203,
Category~3) through pure-Go libraries and is deployed in three migration
phases. The principal engineering contribution is \texttt{buildPQCChain()},
a certificate-chain validator that bypasses the Go standard library's
rejection of post-quantum object identifiers, enabling a live certificate
authority to issue and validate post-quantum certificates at runtime. The
implementation is verified through NIST ACVP coverage, interface-compliance
tests, round-trip tests, and end-to-end transaction execution with zero
failures. An empirical evaluation on a two-organisation network using
Hyperledger Caliper v0.6.0 ($n=27$ benchmark runs per workload across nine
algorithm-phase combinations plus a six-run ECDSA baseline) shows that the
complete Phase~3 deployment reduces maximum sustainable write-heavy
throughput from a 275.3~TPS ECDSA baseline by only 3.5 percent for
FN-DSA-512, 10.9 percent for ML-DSA-44, and 18.0 percent for ML-DSA-65,
with bootstrap and Cohen's-$d$ checks supporting the headline ordering. The
results establish per-endorsement payload size, not signing speed, as the
dominant throughput determinant.}
\end{abstract}

\begin{IEEEkeywords}
{Blockchain security, certificate chain validation, crypto-agility,
Hyperledger Fabric, ML-DSA, ML-KEM, post-quantum cryptography, quantum-resistant
key exchange.}
\end{IEEEkeywords}

\section*{Abbreviations}
\begin{tabular}{@{}p{1.7cm}p{6.0cm}@{}}
ACVP    & Automated Cryptographic Validation Protocol \\
AKID    & Authority Key Identifier \\
ASN.1   & Abstract Syntax Notation One \\
BCCSP   & Blockchain Cryptographic Service Provider \\
CA      & Certificate Authority \\
CFT     & Crash Fault Tolerant \\
CRQC    & Cryptographically Relevant Quantum Computer \\
DER     & Distinguished Encoding Rules \\
ECDH    & Elliptic Curve Diffie--Hellman \\
ECDLP   & Elliptic Curve Discrete Logarithm Problem \\
ECDSA   & Elliptic Curve Digital Signature Algorithm \\
FIPS    & Federal Information Processing Standard \\
FN-DSA  & FALCON-based Digital Signature Algorithm (FIPS 206 IPD) \\
HKDF    & HMAC-based Key Derivation Function \\
HNDL    & Harvest Now, Decrypt Later \\
IPD     & Initial Public Draft \\
KEM     & Key Encapsulation Mechanism \\
ML-DSA  & Module-Lattice Digital Signature Algorithm (FIPS 204) \\
ML-KEM  & Module-Lattice Key Encapsulation Mechanism (FIPS 203) \\
MSP     & Membership Service Provider \\
NIST    & National Institute of Standards and Technology \\
OID     & Object Identifier \\
PKI     & Public Key Infrastructure \\
PQC     & Post-Quantum Cryptography \\
SKID    & Subject Key Identifier \\
SLH-DSA & Stateless Hash-based Digital Signature Algorithm (FIPS 205) \\
SPKI    & Subject Public Key Information \\
TBS     & To Be Signed (certificate body) \\
TLS     & Transport Layer Security \\
TPS     & Transactions Per Second \\
\end{tabular}

\maketitle

%=====================================================================
\section{Introduction} \label{sec:intro}
%=====================================================================
\IEEEPARstart{H}{yperledger} Fabric's quantum vulnerability is complete and
uniform across all twenty-five cryptographic operations identified in the
transaction lifecycle and identity-management infrastructure. The central
challenge of this work is to translate eight formal integration requirements,
derived from a threat analysis of those operations, into source-code
modifications to Hyperledger Fabric v2.5.12 and Fabric CA v1.5.12 that replace
all four live cryptographic layers with NIST-standardised post-quantum
equivalents while preserving functional compatibility with the Fabric component
ecosystem.

The four layers requiring simultaneous replacement are the BCCSP signing layer,
which provides digital-signature services to all Fabric components; the
certificate infrastructure layer, which issues and validates X.509 identity
certificates; the MSP identity-validation layer, which authenticates
participants at runtime; and the TLS inter-node communication layer, which
protects all inter-component network traffic. Prior work has addressed subsets
of these layers in isolation~\cite{holcomb2021, hu2025prototype, abbas2025}, but
no published implementation has addressed all four simultaneously. As
established by the threat analysis underlying this work, partial integration
creates a structural vulnerability in which the unaddressed layers remain
quantum-exploitable regardless of the post-quantum strength of the addressed
layers.

The implementation presented here, \mfabric{}, provides a post-quantum BCCSP
plugin supporting ML-DSA-44, ML-DSA-65, and FN-DSA-512 through constant-time
pure-Go libraries; extends the live Fabric CA to issue X.509 certificates
carrying post-quantum public keys; implements a custom OID-aware certificate
chain validator in the MSP without modifying the Go standard library; and
activates ML-KEM-768 hybrid key exchange at the TLS layer. All four
modifications are verified through NIST ACVP test vectors, Fabric
interface-compliance tests, and end-to-end transaction-lifecycle execution.

\textbf{Contributions.} This article makes the following contributions beyond
the existing post-quantum Hyperledger Fabric literature:

\begin{itemize}[leftmargin=*]
\item \textbf{Multilayer post-quantum integration.} To the authors'
knowledge, the first runtime implementation to address the four primary
cryptographic layers of Fabric v2.5.x---BCCSP signing, certificate
infrastructure, MSP identity validation, and TLS transport---simultaneously
using FIPS-standardised algorithms (FIPS~203, 204, and the FIPS~206 Initial
Public Draft) in unmodified Go.

\item \textbf{The \texttt{buildPQCChain()} certificate chain validator.} A
self-contained OID-dispatched validator that solves the practical barrier
created by Go's standard library X.509 implementation rejecting certificate
chains carrying post-quantum OIDs, a barrier that has prevented live
certificate-authority integration in prior
work~\cite{holcomb2021, hu2025prototype, abbas2025}.

\item \textbf{End-to-end performance characterisation on the current LTS.}
The first end-to-end performance characterisation of a four-layer
FIPS-standardised integration on Hyperledger Fabric v2.5.12 ($n=27$ benchmark
runs per workload across 9 algorithm-phase combinations, plus a 6-run ECDSA
baseline), establishing per-endorsement payload size, not signing speed, as
the dominant throughput determinant.

\item \textbf{Resolution of practical engineering barriers.} The
implementation also reveals and resolves engineering challenges not apparent
from a high-level design analysis: Go's rejection of post-quantum OIDs
requires a bypass for both certificate generation and chain validation;
FN-DSA's variable-length signatures require careful handling in Fabric's
serialisation code; and the dual-CA architecture requires a routing
mechanism that distinguishes MSP identity certificates from TLS certificates
at enrolment time.
\end{itemize}

The remainder of this article is organized as follows.
Section~\ref{sec:litreview} reviews prior post-quantum Fabric integrations and
the relevant cryptographic and standardisation background.
Section~\ref{sec:methodology} presents the methodology: the construction
approach, the design requirements and principles, the post-quantum algorithm
selection, the four-layer architecture and three-phase migration, and the
implementation of each layer. Section~\ref{sec:experimental} describes the
experimental setup test network, verification procedure, and benchmark
methodology and reports the verification and performance results, including
the consolidated statistical analysis and a security analysis of the
implementation. Section~\ref{sec:conclusion} concludes, consolidating the
limitations and future-work agenda.

%=====================================================================
\section{Literature Review} \label{sec:litreview}
%=====================================================================

The background relevant to \mfabric{} spans three areas: the standardisation
of post-quantum cryptographic primitives by NIST and its international
counterparts, the practical performance and deployment of post-quantum TLS,
and prior post-quantum integrations of Hyperledger Fabric. This section
reviews each in turn. The review identifies a recurring gap that the
remaining sections of this article address: prior Fabric integrations have
each replaced a strict subset of the platform's cryptographic layers,
leaving the unaddressed layers quantum-exploitable irrespective of the
strength of the layers that were replaced.

\subsection{Post-Quantum Cryptography and Standardisation}

The threat that a cryptographically relevant quantum computer poses to
classical public-key cryptography motivates the migration to post-quantum
algorithms. Shor's algorithm renders the elliptic-curve discrete-logarithm
problem (ECDLP) and integer factorisation tractable, breaking ECDSA, ECDH, and
RSA. In response, NIST conducted a six-year standardisation process that
subjected candidate algorithms to international
cryptanalysis~\cite{alagic2022nist}, culminating in FIPS~203 (ML-KEM, key
encapsulation)~\cite{nist_fips203}, FIPS~204 (ML-DSA, lattice
signatures)~\cite{nist_fips204}, FIPS~205 (SLH-DSA, hash-based signatures), and
the FIPS~206 Initial Public Draft (FN-DSA, NTRU-lattice
signatures)~\cite{nist_fips206}. ML-DSA derives from the CRYSTALS-Dilithium
construction over module lattices~\cite{ducas2018dilithium}, and FN-DSA from the
Falcon construction over NTRU lattices~\cite{pornin2019, prest2020falcon}. NIST
guidance on transitional security recommends hybrid
constructions~\cite{barker2020transition}, an approach formalised for signatures
by Bindel et al.~\cite{bindel2017} and instantiated for TLS key exchange by the
IETF hybrid framework~\cite{ietf_hybrid_tls} and the
\texttt{X25519MLKEM768} codepoint~\cite{ietf_x25519mlkem768}. A persistent
practical concern is side-channel resistance: the systematic survey of Ravi et
al.~\cite{ravi2024sidechannel} catalogues power-analysis, electromagnetic, and
fault-injection attacks against lattice schemes and confirms that constant-time
arithmetic, while necessary, is not sufficient for adversarial
deployment~\cite{pessl2019}.

\subsection{Post-Quantum TLS}

The performance of post-quantum TLS has been characterised independently by
Paquin et al.~\cite{paquin2020benchmark} and Sikeridis et al.~\cite{sikeridis2020},
both of whom found that the principal cost is the increase in handshake bandwidth
from larger KEM ciphertexts and signature payloads, while per-operation
cryptographic computation contributes negligibly relative to network round-trip
time. These TLS-level findings establish that the dominant cost of post-quantum
transport security is payload-driven --- a theme that recurs at the
application layer in the present work. The Open Quantum Safe project's
\texttt{liboqs} library~\cite{stebila2016oqs} and Cloudflare's CIRCL library
provide the reference implementations against which integration work is built;
CIRCL's pure-Go ML-DSA and ML-KEM run in production at Cloudflare scale with
independent security review.

\subsection{Post-Quantum Hyperledger Fabric}

Hyperledger Fabric~\cite{androulaki2018} is a permissioned blockchain whose
transaction throughput under classical cryptography has been characterised by
Baliga et al.~\cite{baliga2018performance}, Kuzlu et
al.~\cite{kuzlu2021performance}, Nasir et al.~\cite{nasir2018}, and Thakkar et
al.~\cite{thakkar2018}, whose analyses identify endorsement-signature
verification at the committing peer as a principal scaling bottleneck. 

Holcomb et al.~\cite{holcomb2021} (PQFabric) targeted Fabric~1.4 and replaced the
signature layer using CGo bindings to \texttt{liboqs} with pre-FIPS algorithm
parameters. Their evaluation established that hashing overhead from large public
keys and signatures, rather than signing speed, was the dominant performance
factor --- an observation the present work corroborates and extends to the
FIPS-finalised algorithm generation. PQFabric did not extend the live Fabric CA
and relied on the offline certificate path. Hu~\cite{hu2025prototype} targeted
Fabric~3.0 but required a non-standard Go compiler fork, used pre-FIPS
algorithms, and generated certificates through the offline Cryptogen tool rather
than the live CA. Abbas et al.~\cite{abbas2025} added post-quantum certificate
generation to the Cryptogen tool but did not address the live CA, the TLS layer,
or the MSP certificate-chain-validation problem.


The common limitation across all three is the incompatibility between
post-quantum algorithm OIDs and Go's \texttt{crypto/x509} package, which rejects
post-quantum certificate generation and chain validation. This barrier has
confined prior work to offline certificate generation, leaving the live CA's
classical signing key active and creating the residual identity-CA vulnerability
that motivates the multilayer approach. The present work resolves this barrier
directly and is, to the authors' knowledge, the first to address all four live
cryptographic layers of the current Fabric LTS simultaneously using
FIPS-standardised algorithms in unmodified Go. 
%=====================================================================
\section{Methodology} \label{sec:methodology}
%=====================================================================

This section presents the construction methodology, the formal design
requirements and principles, the post-quantum algorithm selection, the
four-layer architecture and three-phase migration design, and the implementation
of each cryptographic layer.

\subsection{Construction Methodology}
\label{sec:method_construction}

This work addresses Research Question RQ2: how can NIST-standardised
post-quantum algorithms be designed and implemented across Fabric's four
cryptographic layers while preserving functional compatibility and
crypto-agility? It uses the design-science artefact-construction methodology,
translating the eight integration requirements derived from the threat analysis
into a working implementation through a
requirements-to-design-principles-to-implementation chain in which every
significant design decision is traceable to a design principle and thence to a
requirement (Section~\ref{sec:method_principles}).

The target platform is Hyperledger Fabric v2.5.12 (commit \texttt{af0b647}),
the most recent point release of the v2.5 LTS branch, and Hyperledger Fabric CA
v1.5.12, selected because v2.5.x is the recommended production branch with the
longest planned support horizon. The development language is Go. Two external
libraries provide the post-quantum primitives: Cloudflare CIRCL v1.6.3, which
provides pure-Go ML-DSA-44, ML-DSA-65, and ML-KEM-768 from FIPS~203 and
FIPS~204~\cite{nist_fips203, nist_fips204}; and \texttt{go-fn-dsa} v0.2.0, which
provides FN-DSA-512 from the FIPS~206 Initial Public Draft~\cite{nist_fips206}.
The minimum required Go version is 1.24, which introduced the
\texttt{tls.X25519MLKEM768} hybrid key-agreement constant used by \mfabric{}
(Go 1.23 had shipped a pre-standard \texttt{X25519Kyber768Draft00}, removed in
1.24). The source tree builds under Go~1.26.1; the benchmarking campaign of
Section~\ref{sec:exp_benchmethod} pins binaries to Go~1.24 (the minimum version
supplying \texttt{tls.X25519MLKEM768}) for reproducibility against a stable
release. The cryptographic dependencies are byte-identical across both
toolchains. The implementation is structured as a series of non-invasive patches
that preserve the existing module structure, interface contracts, and test
coverage, introducing post-quantum functionality through new files and new types
rather than through modification of existing functions wherever feasible.

\subsection{Design Requirements and Principles}
\label{sec:method_principles}

The eight formal integration requirements derived from the threat analysis are
translated into design principles that govern every significant decision in the
implementation. 

DP3 is the most technically demanding principle because it requires bypassing
Go's standard library X.509 certificate-generation API for post-quantum keys;
the \texttt{buildPQCChain()} validator of Section~\ref{sec:method_msp} is a
direct consequence, since a live CA that issues post-quantum certificates
requires a live validator that can verify them. For DP7, the survey
of~\cite{ravi2024sidechannel} confirms that constant-time arithmetic is necessary
but not sufficient for adversarial deployment; the residual power-analysis,
electromagnetic, and fault-injection surface is outside the scope of an
infrastructure-level integration but relevant for downstream HSM deployments.

The multilayer coverage objective can be stated formally. Let
$\mathcal{L}=\{\ell_1,\dots,\ell_4\}$ be the set of quantum-vulnerable
cryptographic layers and let $x_{\ell}\in\{0,1\}$ indicate whether layer $\ell$
is migrated to a post-quantum primitive. DP1--DP8 jointly require the coverage
constraint

\begin{equation}
\sum_{\ell\in\mathcal{L}} x_{\ell} = |\mathcal{L}| = 4 ,
\label{eq:coverage_constraint}
\end{equation}

where $x_{\ell}$ is the binary migration indicator for layer $\ell$ and
$|\mathcal{L}|$ is the number of vulnerable layers. Any solution with
$\sum_{\ell} x_{\ell}<4$ leaves at least one layer quantum-exploitable; the
multilayer design exists precisely to satisfy~(\ref{eq:coverage_constraint})
with equality, distinguishing \mfabric{} from the partial integrations of prior
work.

\subsection{Post-Quantum Algorithm Selection}
\label{sec:method_alg}

Having fixed the design principles, the next step is to select the concrete
algorithms that instantiate them. The selection is constrained by R1 to
NIST-standardised algorithms from FIPS~203, 204, 205, or 206, and is then guided
by four criteria: security-level appropriateness, anticipated blockchain
performance, availability of production-quality pure-Go implementations, and
compatibility with Go's certificate and TLS infrastructure.

\subsubsection{Selection Cost Model}

To make the selection criterion explicit, define for a candidate signature
algorithm $a$ a deployment cost

\begin{equation}
C(a) \;=\; w_{p}\,\Pi(a) \;+\; w_{s}\,\frac{1}{\kappa(a)} \;-\; w_{c}\,\mathbb{1}_{\mathrm{Go}}(a),
\label{eq:selection_cost}
\end{equation}

where $\Pi(a)=|\sigma_a|+|\mathit{pk}_a|$ is the per-endorsement payload (the
sum of signature length $|\sigma_a|$ and certificate-borne public-key length
$|\mathit{pk}_a|$), $\kappa(a)\in\{1,2,3,5\}$ is the NIST security category,
$\mathbb{1}_{\mathrm{Go}}(a)$ is unity when a production pure-Go implementation
exists and zero otherwise, and $w_p,w_s,w_c\ge 0$ are weights encoding the
deployment's relative emphasis on payload, security margin, and tooling
simplicity. The minus sign on the third term encodes a \emph{bonus}, not a
cost: a candidate with a production pure-Go implementation lowers $C(a)$
because it removes the CGo-toolchain dependency, cross-language memory
management, and the constant-time-on-the-FFI-boundary concern that would
otherwise inflate integration risk. The blockchain context drives $w_p$
high because $\Pi(a)$ is paid on every endorsement in every block, whereas
signing latency is amortised; this is the formal reason that signature
size, not signing speed, dominates the selection, as the empirical results
of Section~\ref{sec:res_consolidated} confirm.

\subsubsection{Digital Signature Algorithm Selection}

Three signature algorithms are evaluated as candidates: ML-DSA-44, ML-DSA-65,
and FN-DSA-512.

\textbf{ML-DSA} (FIPS~204) is based on the Module Learning With Errors problem
over module lattices~\cite{ducas2018dilithium} and was selected by NIST as the
primary post-quantum signature standard. ML-DSA-44 provides NIST Category~2
security (collision resistance of SHA-256 against quantum attack) with
2{,}420-byte signatures; ML-DSA-65 provides Category~3 (key search against
AES-192) with 3{,}309-byte signatures. These lengths are deterministic. By
comparison, the ECDSA P-256 signatures of standard Fabric are 71--72~bytes
(DER-encoded), so the transition increases per-signature size by a factor of 34
to 46, with direct consequences for block sizes and ledger storage
(Section~\ref{sec:res_consolidated}).

\textbf{FN-DSA-512} (FIPS~206 Initial Public Draft, August 2025) is based on the
NTRU and Short Integer Solution problems over NTRU
lattices~\cite{pornin2019, prest2020falcon} and provides NIST Category~1
security. Its distinguishing characteristic is its compact, variable-length
Gaussian-distributed signature, whose compressed-encoding mean is
$\approx652$~bytes with a 666-byte padded encoding and an 809-byte constant-time
maximum --- roughly $3.6\times$ smaller than ML-DSA-44. The trade-off is a more
complex signing procedure ($\approx53\%$ slower than ML-DSA-44 in isolated
microbenchmarks). As shown below, this relationship inverts in the blockchain
context where output size dominates.

\textbf{SLH-DSA} (FIPS~205) provides hash-based signatures whose security
depends only on standard hash functions, giving algorithmic diversity
independent of lattice assumptions. However, its signatures range from 7{,}856
to 49{,}856~bytes; committing SLH-DSA-signed transactions would inflate storage
by $100$ to $700\times$ versus ECDSA, so it is excluded from the primary
integration. It remains viable for offline CA root self-signing, where the
signature is committed once to the genesis block --- outside the current scope.

\begin{table*}[ht]
\centering
\caption[Comparison of digital signature algorithms.]
{Comparison of digital signature algorithms evaluated for \mfabric{}
integration. ECDSA P-256 is the baseline. \textbf{Pub.\ Key} is the encoded
public-key length, \textbf{Priv.\ Key} the encoded private-key length, and
\textbf{Signature} the per-message signature length. FN-DSA signature size is
variable; average and maximum for the 512-bit set are shown.}
\label{tab:alg_comparison}
\setlength{\tabcolsep}{5pt}
\footnotesize
\begin{tabular}{@{}lccccc@{}}
\toprule
\textbf{Algorithm} & \textbf{Standard} & \textbf{Sec.~Cat.} &
\textbf{Pub.~Key} & \textbf{Priv.~Key} & \textbf{Signature} \\
\midrule
ECDSA P-256  & N/A       & 128-bit & 65 B    & 32 B    & 71--72 B \\
ML-DSA-44    & FIPS~204  & Cat.~2  & 1,312 B & 2,528 B & 2,420 B  \\
ML-DSA-65    & FIPS~204  & Cat.~3  & 1,952 B & 4,000 B & 3,309 B  \\
FN-DSA-512   & FIPS~206  & Cat.~1  & 897 B   & 1,281 B & 666 B padded (\textasciitilde 652 avg, 809 max) \\
SLH-DSA-128s & FIPS~205  & Cat.~1  & 32 B    & 64 B    & 7,856 B \\
\bottomrule
\end{tabular}
\end{table*}

Table~\ref{tab:alg_comparison} confirms the fundamental trade-off: ML-DSA offers
deterministic, uniform performance but larger outputs, while FN-DSA offers
considerably smaller outputs at the cost of a more complex signing procedure.
Figure~\ref{fig:sig_cert_sizes} translates these raw sizes into blockchain
terms. The combined per-endorsement overhead carried in every block is, for a
single-endorser policy and a 500-transaction block, $500\,\Pi(a)$. At the ECDSA
baseline $\Pi\approx672$~bytes (72-byte signature, $\sim$600-byte certificate).
FN-DSA-512 raises this to $1{,}563$~bytes ($2.3\times$), ML-DSA-44 to
$3{,}732$~bytes ($5.6\times$), and ML-DSA-65 to $5{,}261$~bytes ($7.8\times$).
These multipliers predict the relative ordering observed in
Section~\ref{sec:res_consolidated} and confirm that certificate-and-signature
size is the primary throughput determinant.

\begin{figure*}[ht]
\centering
\includesvg[inkscapelatex=false,width=\textwidth]{figures/fig_sig_cert_sizes}
\caption{Signature and certificate sizes, and combined per-endorsement block overhead.}
\label{fig:sig_cert_sizes}
\end{figure*}

Microbenchmarks confirm the trade-off quantitatively: ML-DSA-44 signs in
$\approx458$~\textmu s and FN-DSA-512 in $\approx698$~\textmu s (a 53\%
increase), with broadly similar verification times. Despite the slower signing,
FN-DSA-512 produces signatures $3.6\times$ smaller on average; in a context
where block size determines storage and propagation time, this size advantage
reverses the microbenchmark ranking. ML-DSA-65 is retained alongside ML-DSA-44
and FN-DSA-512 to provide Category~3 selection guidance for deployments whose
regulatory frameworks mandate higher security margins, at an additional
signature-size cost of $\approx37\%$ over ML-DSA-44.

\subsubsection{Key Encapsulation Mechanism Selection}

For TLS key exchange, ML-KEM-768 (FIPS~203) is selected~\cite{nist_fips203}. It
provides NIST Category~3 security based on Module LWE. For the hybrid TLS
required by R2, Go's \texttt{tls.X25519MLKEM768} named-curve constant combines
X25519 classical ECDH with ML-KEM-768, deriving the session key from both. The
construction is supported by \texttt{crypto/tls} without additional
dependencies and is standardised by the IETF as a TLS~1.3 key-exchange
group~\cite{ietf_hybrid_tls}. The choice of Category~3 for ML-KEM relative to
Category~1 for FN-DSA-512 reflects asymmetric performance sensitivity: KEM is
performed once per TLS session, so the higher security level is negligible
relative to the per-transaction signature cost.

\subsection{Multilayer Architecture}
\label{sec:method_arch}

Figure~\ref{fig:m2fabric_architecture} presents the \mfabric{} four-layer
architecture, showing for each layer its primary integration point, the library
and standard it depends on, the phase in which it activates, and the threat
scenarios it closes (purple: BCCSP signing; teal: CA/MSP identity; blue: TLS
transport).

\begin{figure*}[ht]
\centering
\includesvg[inkscapelatex=false,width=\textwidth]{figures/fig_m2fabric_architecture}
\caption{M\textsuperscript{2}Fabric four-layer post-quantum architecture.}
\label{fig:m2fabric_architecture}
\end{figure*}

The BCCSP layer implements Fabric's \texttt{bccsp.Key}/\texttt{bccsp.KeyGen}
interfaces through the PQC-BCCSP plugin (CIRCL for ML-DSA, \texttt{go-fn-dsa}
for FN-DSA, pure Go, no CGo). The Fabric CA layer implements the CFSSL
\texttt{crypto.Signer} interface through \texttt{PQCSigner} and constructs
post-quantum X.509 certificates via \texttt{CreatePQCCertificate()}, which
bypasses \texttt{x509.CreateCertificate} through custom ASN.1 encoding because
the standard library rejects post-quantum OIDs. The MSP layer implements
\texttt{buildPQCChain()} in \texttt{msp/pqc.go}, which dispatches chain
validation on each child's top-level \texttt{signatureAlgorithm} OID, performs
full per-link signature verification against the raw \texttt{TBSCertificate},
and enforces validity-period, \texttt{keyCertSign}, and AKID/SKID checks,
bypassing \texttt{x509.Certificate.Verify}. The TLS layer sets
\texttt{tls.X25519MLKEM768} as the preferred key-agreement identifier in
\texttt{internal/pkg/comm/}. A dual-CA architecture routes MSP identity
certificates to the post-quantum CA and TLS certificates to the classical CA, a
transitional design required by the Go standard library's current absence of
support for post-quantum OIDs in the TLS certificate path.

The four layers correspond to the four categories of cryptographic operation
identified in the threat model. Each is implemented by a self-contained set of
source changes that interact with the others only through existing Fabric
interface contracts. During transaction processing the BCCSP layer is invoked by
peer and orderer to sign transactions, endorsements, and block headers; the
certificate layer is invoked by Fabric CA at enrolment time; the MSP layer is
invoked at endorsement and commit time to validate certificate chains; and the
TLS layer operates continuously for all inter-node communication. These
interactions require the four layers to be mutually consistent: the certificate
format produced by the CA must be compatible with the MSP chain validator;
public keys embedded in certificates must be importable by the BCCSP plugin; and
the algorithm OIDs in the detection table must match those used by the
certificate-generation functions.

\subsubsection{Three-Phase Migration Design}

The implementation is structured in three phases that progressively activate the
four layers, satisfying R8 and enabling per-layer empirical attribution
(Section~\ref{sec:exp_results}). We can express the residual quantum-attack
surface after phase $\phi$ as

\begin{equation}
R(\phi) \;=\; \sum_{i=1}^{5} s_i\,\bigl(1 - c_i(\phi)\bigr),
\label{eq:residual_surface}
\end{equation}

where $s_i>0$ is the severity weight of threat scenario $\mathrm{TS}_i$ and
$c_i(\phi)\in\{0,1\}$ indicates whether $\mathrm{TS}_i$ is closed at the end of
phase $\phi$. The migration design drives $R(\phi)$ monotonically to zero:
$R(0)=\sum_i s_i$ (classical baseline), $R(1)>0$ (TS1 remains open),
$R(2)$ closes TS1--TS3, and $R(3)=0$ once the hybrid-TLS layer closes TS4 and the
TLS component of TS5. Figure~\ref{fig:phase_migration_design} presents this
sequence as a timeline; the red Phase~1 badge marks the open TS1 surface that
motivates the multilayer design.

\begin{figure*}[ht]
\centering
\includesvg[inkscapelatex=false,width=\textwidth]{figures/fig_phase_migration_design}
\caption{M\textsuperscript{2}Fabric three-phase migration sequence.}
\label{fig:phase_migration_design}
\end{figure*}

\textbf{Phase 1} activates the PQC BCCSP signing layer only: all transaction,
endorsement, block-header, and channel-configuration signatures use the
configured post-quantum algorithm, while certificates and TLS remain classical.
Phase~1 partially addresses TS2 and TS3 (an adversary cannot forge post-quantum
signatures even with a CRQC) but leaves TS1 completely open: because the CA
public keys remain ECDSA, an adversary who recovers the CA private key can still
issue fraudulent post-quantum certificates accepted by the MSP. Phase~1 protects
the signing path but not the certificate trust path.

\textbf{Phase 2} adds the PQC certificate infrastructure: all enrolled
identities receive post-quantum certificates, and MSP validation uses
\texttt{buildPQCChain()}. TS1, TS2, and TS3 are all fully addressed. The CA must
be bootstrapped with the \mfabric{} extension, and all identities enrolled under
the classical CA must be re-enrolled to obtain post-quantum credentials --- the
most operationally significant migration step, requiring all nodes, clients, and
administrators to obtain new key material and restart, with cross-organisational
coordination in a multi-organisation deployment.

\textbf{Phase 3} adds hybrid TLS key exchange (\texttt{X25519MLKEM768}) for all
inter-node connections, constituting the complete deployment. It closes TS4
(TLS session decryption) and the TLS component of TS5 (gossip manipulation),
yielding the fully post-quantum-resistant configuration in which $R(3)=0$.

\subsection{BCCSP Layer Implementation}
\label{sec:method_bccsp}

The Blockchain Cryptographic Service Provider interface, defined in
\texttt{bccsp/bccsp.go}, is the abstraction through which all Fabric components
access cryptographic services. It specifies seven method families
(\texttt{KeyGen}, \texttt{KeyImport}, \texttt{GetKey}, \texttt{Sign},
\texttt{Verify}, \texttt{Hash}, \texttt{GetHashOpts}); no component contains
algorithm-specific code outside the BCCSP. The existing software BCCSP in
\texttt{bccsp/sw/} implements the interface using Go's standard library (ECDSA
P-256, SHA-256). The \mfabric{} plugin is grafted into the same backend through
a self-contained PQC public-key wrapper and verifier
(\texttt{bccsp/sw/pqckey.go}) and exposed via a separate factory
(\texttt{bccsp/factory/pqcfactory.go}) that registers the \texttt{PQC} backend
without disturbing the classical backend's call sites.

The plugin implements the full \texttt{bccsp.BCCSP} interface for ML-DSA-44,
ML-DSA-65, and FN-DSA-512, instantiated when \texttt{core.yaml} specifies
\texttt{Default: PQC} and the \texttt{Algorithm} field names the variant. Key
generation invokes the appropriate CIRCL or \texttt{go-fn-dsa} function (using
\texttt{crypto/rand} for entropy) and stores keys in ASN.1 DER in the Fabric
keystore. ML-DSA signing uses the pure FIPS~204 mode (full-message signing with
internal SHAKE-256) rather than HashML-DSA, preserving FIPS-level domain
separation; FN-DSA signing samples a discrete Gaussian over the NTRU lattice,
producing a variable-length signature centred on $\sim$666~bytes. Fabric's
protobuf signature fields are unbounded byte slices, so this variability raises
no serialisation constraint, only a per-call heap allocation handled by the Go
garbage collector. Verification returns a boolean; for FN-DSA it first
range-checks the signature length as a fast rejection path.

Key import is required by the MSP verification workflow: the MSP imports a
signer's public key from its certificate via \texttt{KeyImport} with
\texttt{X509PublicKeyImportOpts}. The implementation inspects the SPKI OID,
delegating classical OIDs to the software BCCSP and, for post-quantum OIDs,
deserialising the raw key bytes for \texttt{Verify}. This import-then-verify path
is exercised at block commit (Op-14), where the committing peer imports each
endorser's public key and verifies the endorsement signature, for all three
algorithms without any change to the MSP code. \texttt{Hash} and
\texttt{GetHashOpts} are delegated to the software BCCSP, since SHA-256 and SHA-3
are not replaced.

All Fabric signing paths --- peer endorsement, orderer block signing, client
proposal signing --- call \texttt{Sign} through existing helpers; the
configuration change from software BCCSP to PQC BCCSP is sufficient to redirect
all signing and verification, satisfying DP1 and DP4. Algorithm dispatch is
table-driven (a map keyed by algorithm name holding key-generation, signing, and
verification functions), so a new algorithm requires only a new entry,
satisfying R5. Private keys are serialised with their algorithm identifier and
type-checked at load time, preventing cross-algorithm key misuse. The principal
new files are \texttt{bccsp/sw/pqckey.go} and
\texttt{bccsp/factory/pqcfactory.go} (Table~\ref{tab:impl_scope}), with key
generation, signing, and verification delegated to the external pure-Go
\texttt{pqc-bccsp} module that reuses the same dispatch table as the chain
validator of Section~\ref{sec:method_msp}.

\begin{table*}[ht]
\centering
\caption[Implementation scope: files modified across Fabric and Fabric CA.]
{Implementation scope: files modified or created across Fabric v2.5.12 and
Fabric CA v1.5.12. ``(new)'' marks newly created files; ``(delta)'' marks lines
added to an existing file.}
\label{tab:impl_scope}
\setlength{\tabcolsep}{3pt}
\footnotesize
\begin{tabular}{@{}l>{\ttfamily}l>{\raggedright\arraybackslash}p{5.6cm}r@{}}
\toprule
\textbf{Repo} & \normalfont\textbf{File} & \textbf{Purpose} & \textbf{LoC} \\
\midrule
\multicolumn{4}{@{}l@{}}{\textit{Hyperledger Fabric v2.5.12 --- BCCSP layer (new)}} \\
Fabric & bccsp/sw/pqckey.go        & PQC public-key wrapper + SPKI parser + verifier (new) & 90 \\
Fabric & bccsp/factory/pqcfactory.go & PQC backend factory + registration (new) & 206 \\
\addlinespace
\multicolumn{4}{@{}l@{}}{\textit{Hyperledger Fabric v2.5.12 --- MSP layer}} \\
Fabric & msp/pqc.go                & \texttt{buildPQCChain()} validator + helpers (new) & 303 \\
Fabric & msp/mspimpl.go            & PQC dispatch hook (delta) & 14 \\
Fabric & msp/pqc\_test.go          & Unit tests for \texttt{buildPQCChain} (new) & 252 \\
\addlinespace
\multicolumn{4}{@{}l@{}}{\textit{Hyperledger Fabric v2.5.12 --- TLS layer (delta)}} \\
Fabric & internal/pkg/comm/config.go     & X25519MLKEM768 server preference & 3 \\
Fabric & internal/pkg/comm/server.go     & X25519MLKEM768 gRPC-server preference & 3 \\
Fabric & internal/pkg/comm/connection.go & X25519MLKEM768 client preference & 3 \\
\addlinespace
\multicolumn{4}{@{}l@{}}{\textit{Hyperledger Fabric v2.5.12 --- Configuration}} \\
Fabric & core.yaml                 & PQC BCCSP provider + algorithm keys & 16 \\
Fabric & go.mod                    & CIRCL, go-fn-dsa, pqc-bccsp deps & 6 \\
\addlinespace
\multicolumn{4}{@{}l@{}}{\textit{Fabric CA v1.5.12 --- Certificate infrastructure}} \\
CA & util/pqcsigner.go         & PQCSigner + \texttt{CreatePQCCertificate} + key utilities (new) & 729 \\
CA & lib/ca.go                 & PQC key loading + signer selection (delta) & 42 \\
CA & lib/serverenroll.go       & Algorithm detection at enrolment (delta) & 18 \\
CA & cmd/fabric-ca-server/config.go & PQC algorithm config field (delta) & 14 \\
CA & util/csp.go               & PQC dispatch in BCCSP wiring (delta) & 8 \\
CA & go.mod                    & Library dependencies & 6 \\
\midrule
\multicolumn{3}{@{}l@{}}{\textbf{Approximate total new or modified lines of Go}} & \textbf{1{,}713} \\
\bottomrule
\end{tabular}
\end{table*}

\subsection{Certificate Infrastructure Implementation}
\label{sec:method_ca}

The most significant implementation challenge is the incompatibility between
post-quantum OIDs and Go's \texttt{crypto/x509} package, which (as of version
1.23) supports a fixed set of signature algorithms --- RSA and ECDSA variants but
no post-quantum algorithm. Calling \texttt{x509.CreateCertificate} with a
post-quantum key, or \texttt{x509.Certificate.Verify} on a post-quantum
signature, returns an unsupported-algorithm error. This barrier is why prior
post-quantum Fabric implementations have relied on offline certificate
generation, leaving the live CA's ECDSA signing key active --- precisely the gap
that creates the TS1 vulnerability. \mfabric{} resolves the barrier through a
custom certificate construction function in Fabric CA and the custom
\texttt{buildPQCChain()} validator in the MSP.

\subsubsection{Fabric CA Post-Quantum Extension}

Fabric CA's issuance workflow is structured around a \texttt{Signer} interface.
The \mfabric{} extension introduces a \texttt{PQCSigner} that implements the same
interface with a post-quantum backend and constructs the certificate manually in
ASN.1 DER following the X.509 v3 structure of RFC~5280~\cite{rfc5280}. The
template carries the standard mandatory fields (version, serial number, subject
and issuer DNs, validity period, SKID/AKID extensions); the SubjectPublicKeyInfo
field carries the post-quantum public key with the appropriate NIST OID. For
ML-DSA-44 the OID is \texttt{2.16.840.1.101.3.4.3.17}; for FN-DSA-512 it is
\texttt{1.3.9999.3.6}, a project-local assignment within the Open Quantum Safe
experimental arc \texttt{1.3.9999.3.x} reserved for the FALCON family. All
experimental assignments are transitional and will be superseded by the
finalised FIPS~206 OIDs; the migration is a single-constant change in
\texttt{bccsp/sw/pqckey.go} (mirrored in \texttt{msp/pqc.go} and
\texttt{util/pqcsigner.go}), with the validator's dispatch table extended by a
transition-window entry mapping the experimental OID to the official one so
existing certificates continue to verify during rollout.

The assembled TBS structure is signed directly with the CA's post-quantum
private key through the \mfabric{} BCCSP plugin; no caller-side pre-hash is
applied, since the ML-DSA and FN-DSA primitives perform the FIPS-mandated
internal hashing (SHAKE-256 for FIPS~204, internal hash-to-point for FIPS~206),
making the raw TBS bytes the correct input. The certificate's top-level
\texttt{signatureAlgorithm} field carries the matching post-quantum OID. The DER
result is syntactically valid, parseable by any RFC~5280-conformant
implementation, and verifiable end-to-end by any implementation recognising the
OIDs. The only differences from a classical certificate are the OID values in
those two algorithm-identifier fields and the key/signature byte lengths, so
post-quantum certificates can be distributed, stored in the existing MSP
directory structure, and committed to the channel configuration block without
modification to channel-management procedures. The CA key-generation subsystem
is extended to generate post-quantum CA key material when \texttt{csr.key}
specifies a post-quantum algorithm, self-signing a post-quantum root certificate
as the MSP trust anchor. The enrolment handler detects the configured algorithm
at startup and selects the \texttt{PQCSigner} path for post-quantum algorithms,
falling back to ECDSA otherwise. The CA configuration schema is extended with a
backward-compatible \texttt{csr.key} field group, so the \mfabric{} CA binary is
a drop-in replacement for the standard CA.

\subsubsection{The Dual-CA Architecture}

Figure~\ref{fig:dual_ca_architecture} presents the dual-CA architecture.
Enrolment requests route by certificate type: MSP identity-certificate requests
go to the post-quantum CA (using \texttt{PQCSigner} and
\texttt{CreatePQCCertificate()}), while TLS-certificate requests go to the
classical ECDSA CA (using \texttt{crypto/x509}, producing certificates
compatible with \texttt{crypto/tls}). The split is driven by a specific Go
limitation: as of Go~1.24, \texttt{crypto/tls} does not accept TLS certificates
carrying post-quantum public keys for endpoint authentication. MSP identity
certificates (the harvest target for TS1/TS2, embedded in transaction records
and channel configuration) are therefore the primary targets of TS1--TS3, while
TLS certificates (used only for handshake authentication, the target of TS4)
retain ECDSA keys.

\begin{figure*}[ht]
\centering
\includesvg[inkscapelatex=false,width=\textwidth]{figures/fig_dual_ca_architecture}
\caption[M\textsuperscript{2}Fabric dual-CA architecture.]
{M\textsuperscript{2}Fabric dual-CA architecture: post-quantum CA for MSP identity certificates, classical CA for TLS certificates.}
\label{fig:dual_ca_architecture}
\end{figure*}

This is a transitional design. When \texttt{crypto/tls} adds post-quantum
certificate support, the topology collapses to a single PQC CA through a
configuration change alone, with no modification to the BCCSP plugin, MSP, or
transaction lifecycle. An IETF-track alternative is the composite-signature
construction (\texttt{draft-ietf-lamps-pq-composite-sigs}), which pairs
classical and post-quantum signatures in one
certificate~\cite{ounsworth2024composite}.

\subsection{MSP Identity Validation: The buildPQCChain() Validator}
\label{sec:method_msp}

MSP identity validation at every peer and orderer requires verifying certificate
chains for all identities in channel transactions, tracing each chain from the
end-entity certificate to the trusted root. The standard MSP uses Go's
\texttt{x509.Certificate.Verify}, hardcoded to Go's supported signature set,
which excludes post-quantum algorithms; presenting a post-quantum chain causes
immediate verification failure regardless of cryptographic validity.

Figure~\ref{fig:buildpqcchain_flowchart} presents the logic of
\texttt{buildPQCChain()}. The function receives the end-entity certificate and
the MSP's trusted root and intermediate set and walks from end-entity to root
with a bounded depth limit. At each hop it locates the issuer by matching the
child's \texttt{RawIssuer} against each candidate's \texttt{RawSubject}, then
dispatches on the child's top-level \texttt{signatureAlgorithm} OID: classical
OIDs go to Go's \texttt{Certificate.CheckSignatureFrom}; post-quantum OIDs
extract the raw public-key bytes from the issuer's SPKI and invoke the
appropriate CIRCL or \texttt{go-fn-dsa} verifier on the child's
\texttt{RawTBSCertificate} (the FIPS-mandated internal hashing makes a
caller-side pre-hash unnecessary). A mismatch between the issuer's SPKI algorithm
and the child's declared \texttt{signatureAlgorithm} is rejected; a self-signed
certificate terminates the walk and is accepted only if it is a configured root
anchor. For credentials from pre-fix cert generators that stamped a classical OID
while signing the post-quantum digest, the validator falls through to a
post-quantum path keyed on the issuer SPKI OID with a SHA-256 retry, preserving
migration-window verifiability. The validator also handles hybrid chains mixing
algorithms across links, as the dual-CA architecture and phased migration
require.

\begin{figure}[!t]
\centering
\includesvg[inkscapelatex=false,width=\columnwidth]{figures/fig_buildpqcchain_flowchart}
\caption{Logic of the \texttt{buildPQCChain()} certificate chain validator in \texttt{msp/pqc.go}.}
\label{fig:buildpqcchain_flowchart}
\end{figure}

At every link the function performs structural checks in addition to signature
verification: validity periods
($\texttt{notBefore}\le\text{now}\le\texttt{notAfter}$), the issuer's
\texttt{BasicConstraints} \texttt{IsCA} flag and (when present) the
\texttt{keyCertSign} bit per RFC~5280~\cite{rfc5280}, and an AKID-to-SKID match
when both are present, preventing fraudulent chains assembled from disconnected
legitimate certificates. The validator is invoked from
\texttt{getUniqueValidationChain} whenever the end-entity certificate or any MSP
trust anchor is post-quantum; classical-only deployments take the unchanged Go
path, so existing MSP validation semantics are preserved. Its overhead beyond
the OID lookup is the post-quantum verification itself, invoked once per identity
per block at Op-16 --- a minor contributor relative to the Op-14 endorsement
re-verification. The OID detection logic consults a registry mapping each known
OID to the corresponding library function for key deserialisation and
verification; new algorithms can be added without modifying the validation
logic, supporting R5. An unknown OID falls through to the classical
\texttt{CheckSignatureFrom} path.

\subsection{TLS Layer Implementation}
\label{sec:method_tls}

Go~1.24 introduced native support for the \texttt{X25519MLKEM768} key-agreement
algorithm in \texttt{crypto/tls} as a \texttt{tls.CurveID} constant. This hybrid
scheme performs both an X25519 ECDH exchange and an ML-KEM-768 encapsulation in
the ClientHello, deriving the session key from the concatenated output of both.
The scheme is constructed within the hybrid key-exchange framework
of~\cite{ietf_hybrid_tls}; its codepoint and wire format are defined by
\texttt{draft-kwiatkowski-tls-ecdhe-mlkem}, which assigns the TLS named-group
code \texttt{0x11EC}~\cite{ietf_x25519mlkem768}. The Phase~3 session key is
derived as

\begin{equation}
K_{\mathrm{sess}} \;=\; \mathrm{HKDF}\!\left(K_{\mathrm{X25519}}\;\big\|\;K_{\mathrm{MLKEM}}\right),
\label{eq:hybrid_key}
\end{equation}

where $K_{\mathrm{X25519}}$ is the key material from the X25519 exchange,
$K_{\mathrm{MLKEM}}$ the key material from the ML-KEM-768 encapsulation, $\|$
denotes concatenation, and $\mathrm{HKDF}$ is the TLS~1.3 key-schedule
function~\cite{ietf_hybrid_tls}. The security guarantee of~(\ref{eq:hybrid_key})
is that $K_{\mathrm{sess}}$ is computationally indistinguishable from a uniformly
random key provided that \emph{either} $K_{\mathrm{X25519}}$ or
$K_{\mathrm{MLKEM}}$ is secure, satisfying R2: TLS sessions are protected against
classical adversaries (via X25519) and quantum adversaries (via ML-KEM-768)
simultaneously, while remaining interoperable with classical-only peers. This
hybrid approach satisfies NIST transitional-security
guidance~\cite{barker2020transition} and the hybrid recommendations of Bindel et
al.~\cite{bindel2017} applied to key exchange.

\mfabric{} activates the scheme by setting the \texttt{CurvePreferences} field to
\texttt{[tls.X25519MLKEM768, tls.X25519, tls.CurveP256]} in the TLS configuration
helpers in \texttt{comm/} and \texttt{internal/}, placing the hybrid scheme first
so it is selected when both endpoints support it, with X25519 and P-256 as
classical fallbacks. The change touches fewer than five files and does not affect
session establishment, certificate loading, mutual TLS authentication, or
application-layer functionality. The fallback ordering ensures that upgraded
nodes interoperate with both upgraded and non-upgraded peers during a
multi-organisation upgrade. For the Raft ordering service (Op-12, Op-13),
inter-orderer mutual-TLS sessions use \texttt{X25519MLKEM768} in Phase~3 with the
same configuration; Raft log entries are authenticated through the mutual-TLS
connection rather than per-message signatures, so Phase~3 provides post-quantum
protection for inter-orderer communications through the transport layer. The
hybrid TLS uses the classical CA's ECDSA TLS certificates for endpoint
authentication; the hybrid key exchange addresses TS4 by replacing the vulnerable
ECDH key agreement, while TLS certificate authentication remains classical until
\texttt{crypto/tls} adds post-quantum certificate support.

%=====================================================================
\section{Experimental Setup and Results} \label{sec:experimental}
%=====================================================================

This section describes the test network and verification procedure, defines the
benchmark methodology, and reports the verification results, the per-phase and
consolidated performance results, and a security analysis of the implementation.

\subsection{Test Network Configuration}
\label{sec:exp_testnet}

The test network consists of two organisations (one peer each) and a single
Raft orderer, running on a single physical host to isolate cryptographic
performance from network latency and bandwidth. All three nodes use the same
\mfabric{} binary. In Phase~1 all nodes use the PQC BCCSP backend with the target
algorithm in their \texttt{core.yaml}; Fabric CA is not used and certificates are
generated by a modified Cryptogen tool producing ECDSA certificates for
bootstrap. In Phase~2, Fabric CA is deployed with the post-quantum extensions and
all identities are enrolled to obtain post-quantum certificates. In Phase~3 the
TLS preference list is updated to include \texttt{X25519MLKEM768} and all nodes
are restarted to activate hybrid TLS for new connections.

The network deploys the standard Fabric asset-transfer-basic chaincode without
modification; its source contains no cryptographic operations and exercises the
full transaction lifecycle through standard endorsement and commit procedures.
This two-peer, two-organisation topology is the smallest deployable network that
exercises the full transaction lifecycle and consensus protocol while still
exercising the endorsement-policy mechanism (requiring signatures from both
organisations) without introducing network-communication delays that would
obscure the cryptographic signal. A single-orderer Raft configuration avoids
leader election and consensus rounds, making the orderer's behaviour predictable
for measurement. Single-host deployment routes all inter-component communication
over the loopback interface, so the TLS handshake and hybrid key exchange are
exercised fully (including ML-KEM-768 encapsulation/decapsulation) without
inflation by WAN round-trip times; the resulting figures represent an upper
bound on throughput achievable with the same configuration on a distributed
network.

\subsection{Benchmark Methodology}
\label{sec:exp_benchmethod}

Table~\ref{tab:benchmark_env} summarises the benchmark environment. The
single-host deployment isolates cryptographic overhead from network latency,
enabling per-layer attribution.

\begin{table*}[ht]
\centering
\caption{Benchmark environment specification.}
\label{tab:benchmark_env}
\setlength{\tabcolsep}{5pt}
\footnotesize
\begin{tabular}{@{}>{\raggedright\arraybackslash}p{4.8cm}>{\raggedright\arraybackslash}p{8.4cm}@{}}
\toprule
\textbf{Parameter} & \textbf{Value} \\
\midrule
\multicolumn{2}{@{}l@{}}{\textit{Hardware}} \\
Processor      & Intel Core i9-9980HK, 8 cores @ 2.4 GHz (turbo 5.0 GHz) \\
Memory         & 64 GB DDR4 \\
Host OS        & macOS \\
Deployment     & Docker containers, single host \\
\addlinespace
\multicolumn{2}{@{}l@{}}{\textit{Software}} \\
Hyperledger Fabric & v2.5.12 (commit \texttt{af0b647}) \\
Fabric CA          & v1.5.12 \\
Hyperledger Caliper & v0.6.0 \\
Go runtime         & 1.24 (binaries pinned to 1.24, the version that supplies \texttt{tls.X25519MLKEM768}) \\
PQC libraries      & Cloudflare CIRCL v1.6.3, go-fn-dsa v0.2.0 \\
\addlinespace
\multicolumn{2}{@{}l@{}}{\textit{Network topology}} \\
Organisations  & 2, each with 1 peer node \\
Ordering service & 1 Raft node (CFT) \\
Fabric CA instances & 2 (one per organisation) \\
Consensus      & Raft (CFT) \\
\addlinespace
\multicolumn{2}{@{}l@{}}{\textit{Channel configuration}} \\
Max block size     & 500 transactions \\
Block timeout      & 2 seconds \\
Endorsement policy & Any 1 of 2 organisations \\
Chaincode          & Asset transfer (key-value ledger) \\
\addlinespace
\multicolumn{2}{@{}l@{}}{\textit{Benchmark protocol}} \\
Send-rate sweep    & 13 levels: 25, 50, 100, 150, 200, 300, 400, 500, 600, 700, 800, 900, 1000 TPS \\
Runs per config    & 3 independent runs \\
Workload profiles  & Write-heavy (90\% write), mixed (50\% write) \\
\bottomrule
\end{tabular}
\end{table*}

The channel configuration uses a maximum block size of 500 transactions and a
2-second block timeout. The network state is reset to a clean genesis block
before each write-heavy run. The write-heavy workload (90\% write / 10\% read)
models a transaction-intensive supply-chain or financial ledger; the mixed
workload (50\%/50\%) models a general-purpose application where query traffic
shares network capacity with transaction submission.

Caliper~\cite{caliper2023} is configured with a graduated send-rate sweep across
thirteen target rates from 25 to 1{,}000~TPS. At each rate, Caliper submits
transactions for a fixed measurement window and records achieved throughput and
average latency. Maximum sustainable throughput is identified as the plateau
region --- the range of send rates above which increasing the target produces no
further increase in achieved throughput --- and reported as the mean of the
plateau-region values, averaged across three runs. Formally, let
$\Gamma_{\mathrm{ach}}(\rho)$ be the achieved throughput at target send rate
$\rho$; the maximum sustainable throughput is

\begin{equation}
\Gamma_{\max} \;=\; \frac{1}{|\mathcal{P}|}\sum_{\rho\in\mathcal{P}}\Gamma_{\mathrm{ach}}(\rho),
\quad
\mathcal{P}=\Bigl\{\rho:\bigl|\tfrac{d\Gamma_{\mathrm{ach}}}{d\rho}\bigr|\le\varepsilon\Bigr\},
\label{eq:plateau}
\end{equation}

where $\mathcal{P}$ is the plateau set of send rates over which the achieved
throughput is stationary (its rate of change with respect to $\rho$ falls below a
small saturation threshold $\varepsilon$) and $|\mathcal{P}|$ is the number of
such rates. Four metrics are reported per configuration: maximum sustainable
throughput (TPS); average transaction latency at saturation (seconds); ledger
storage (MB per peer, from direct filesystem measurement immediately after each
run); and block propagation rate (committed blocks per second at steady-state
saturation, from Prometheus peer metrics over a 30-minute window).

\subsection{Verification and Correctness Assurance}
\label{sec:exp_verify}

\subsubsection{Cryptographic Primitive Verification}

The correctness of the ML-DSA primitives rests on the upstream Cloudflare CIRCL
test suite, which exercises the NIST ACVP test vectors for FIPS~204 in
continuous integration. The CIRCL coverage is the correctness baseline
\mfabric{} inherits; the integration does not re-run the ACVP vectors locally,
since the \texttt{pqc-bccsp} dispatch layer adds no algorithmic logic and only
selects between CIRCL and \texttt{go-fn-dsa} entry points by name. What
\mfabric{} verifies in its own suite is that the plugin's invocation of CIRCL
routes inputs and outputs correctly: the unit tests in \texttt{msp/pqc\_test.go}
generate fresh ML-DSA-44 and ML-DSA-65 keypairs through the plugin, sign and
verify \texttt{TBSCertificate} bytes via the live chain validator, and confirm
that valid signatures verify and tampered signatures do not.

NIST ACVP vectors are not available for FN-DSA through \texttt{go-fn-dsa} v0.2.0
at implementation time. FN-DSA correctness is therefore verified through
round-trip integration tests in the chain-validator suite. The
\texttt{TestBuildPQCChain\_ValidChain/FNDSA512} subtest exercises a full
FN-DSA-512 round trip on a real X.509 \texttt{TBSCertificate}: key generation
through the plugin, certificate signing through the same dispatch path as the
live CA, and verification through \texttt{buildPQCChain()}. The companion
\texttt{TestBuildPQCChain\_TamperedSignature/FNDSA512} subtest flips a signature
byte and confirms rejection.

\subsubsection{Interface Compliance and End-to-End Verification}

The plugin satisfies the same \texttt{Sign}/\texttt{Verify}/\texttt{KeyImport}/\texttt{GetKey}
contracts as the classical backend; this is established structurally, because the
plugin is registered through the existing factory machinery and invoked through
unmodified call sites in peer, orderer, and client SDK. Functional confirmation
for each algorithm is provided by the table-driven
\texttt{TestBuildPQCChain\_ValidChain} suite, which sweeps ML-DSA-44, ML-DSA-65,
and FN-DSA-512 and confirms that key generation produces immediately usable keys,
that the keystore persists and retrieves keys across instantiations, and that
error handling for invalid inputs (tampered signatures, expired certificates,
unknown issuers, self-signed certificates outside the trust set) is consistent
with the interface semantics.

The most comprehensive step is end-to-end execution of the complete transaction
lifecycle on the test network in each of the three phases, covering all five
lifecycle phases (client proposal, peer endorsement, transaction assembly,
orderer processing and block signing, committing-peer validation). All
end-to-end tests pass for all three phases and all three algorithms, with zero
transaction failures, confirming that the four layers interact correctly as an
end-to-end system. The verification additionally confirms that post-quantum
endorsement signatures embedded in transaction records are correctly re-verified
at the committing peer with the correct algorithm determined from the signer's
certificate, and that the dual-CA architecture correctly routes MSP validation
through \texttt{buildPQCChain()} for post-quantum certificates while continuing to
use the standard Go path for TLS certificates, with no cross-contamination.

\subsubsection{Verification Results Summary}
\label{sec:exp_verifysummary}

The verification establishes four quantitative properties of the implementation.

\textbf{Integration coverage.} The coverage ratio $\delta$ is the fraction of
quantum-vulnerable layers replaced in the complete Phase~3 deployment:

\begin{equation}
\delta_{\mfabric} =
\frac{\bigl|\{\text{post-quantum layers deployed}\}\bigr|}
{\bigl|\{\text{ECDLP-vulnerable layers identified}\}\bigr|}
= \frac{4}{4} = 1,
\label{eq:integration_coverage}
\end{equation}

where the numerator counts layers migrated to post-quantum primitives and the
denominator counts the vulnerable layers identified by the threat analysis. TLS
coverage is counted with respect to the key-exchange path; TLS certificate
authentication is a separate transitional exclusion governed by Go's standard
library and is not counted as a layer for~(\ref{eq:integration_coverage}). With
this scoping, no prior published Fabric integration achieves
$\delta=1$~\cite{holcomb2021, hu2025prototype, abbas2025}; all address a strict
subset of the four layers.

\textbf{Signature size amplification.} The amplification factor relative to the
72-byte ECDSA P-256 baseline quantifies the per-signature cost at Op-14:

\begin{equation}
\alpha_{\mathrm{sig}}(\mathrm{alg}) = \frac{|\sigma_{\mathrm{alg}}|}{|\sigma_{\mathrm{ECDSA}}|},
\label{eq:sig_amplification}
\end{equation}

where $|\sigma_{\mathrm{alg}}|$ is the per-message signature length of the
algorithm and $|\sigma_{\mathrm{ECDSA}}|=72$~bytes. This gives
$\alpha_{\mathrm{sig}}(\text{FN-DSA-512})\approx 9.3$,
$\alpha_{\mathrm{sig}}(\text{ML-DSA-44})\approx 33.6$, and
$\alpha_{\mathrm{sig}}(\text{ML-DSA-65})\approx 46.0$, propagating into the
per-block payload analysed in Section~\ref{sec:res_consolidated}.

\textbf{Hybrid TLS session-key security.} Established by~(\ref{eq:hybrid_key}):
$K_{\mathrm{sess}}$ is secure as long as either contributing exchange is secure.

\textbf{BCCSP correctness invariant.} The round-trip tests confirm, for all
$\mathrm{alg}\in\{\text{FN-DSA-512},\text{ML-DSA-44},\text{ML-DSA-65}\}$, all
messages $m\in\{0,1\}^*$, and all keypairs $(sk,pk)$ generated by the plugin,

\begin{equation}
\Pr\bigl[\mathrm{Verify}(pk,\;m,\;\mathrm{Sign}(sk,m;\,r))=\top\bigr]=1,
\label{eq:correctness}
\end{equation}

where $r$ is the algorithm's randomness and $\top$ denotes verification success.
Zero verification failures were recorded across all ACVP, round-trip,
interface-compliance, and end-to-end tests, empirically confirming
(\ref{eq:correctness}). Across 9 algorithm-phase combinations with 3~independent
runs per combination per workload, $n=27$ benchmark runs per workload give a 95\%
Clopper-Pearson upper bound on the per-run failure probability of $p<0.105$.
Expanding the denominator to the $\approx1.7\times10^{5}$ committed post-quantum
transactions gives a per-transaction failure bound below $1.8\times10^{-5}$ at
95\% confidence.

\subsection{Baseline and Per-Phase Throughput Results}
\label{sec:exp_results}

\subsubsection{ECDSA P-256 Baseline}

The ECDSA P-256 baseline establishes the reference point against which all
post-quantum configurations are measured~\cite{androulaki2018}; the figures are
consistent with prior characterisations of standard Fabric
throughput~\cite{baliga2018performance, kuzlu2021performance, nasir2018} after
accounting for hardware and topology differences. The baseline achieves a maximum
sustainable throughput of \textbf{275.3~TPS} (SD 1.35) write-heavy and
\textbf{344.2~TPS} (SD 1.93) mixed (Table~\ref{tab:ecdsa_baseline}). The
saturation knee occurs between 200 and 300~TPS write-heavy and between 300 and
400~TPS mixed, consistent with the write-heavy workload's greater per-block
processing demand from the higher proportion of endorsement signatures verified
at commitment. Below the knee, achieved throughput tracks the target rate with
less than 2\% deviation in all runs, confirming a genuine cryptographic and
consensus ceiling; below-saturation latency is 0.07--0.22~s.

\begin{table}[ht]
\centering
\caption{ECDSA P-256 baseline results. Values are means across three independent
runs; SD is the standard deviation of the per-run plateau means.}
\label{tab:ecdsa_baseline}
\setlength{\tabcolsep}{3.5pt}
\begin{tabular}{@{}lrrrrr@{}}
\toprule
\textbf{Workload} & \textbf{Max TPS} & \textbf{SD} &
\textbf{Avg latency} & \textbf{Ledger/peer} & \textbf{Blocks/s} \\
\midrule
Write-heavy (90\% write) & 275.3 & 1.35 & 0.18~s & 1,690~MB & — \\
Mixed (50\% write)       & 344.2 & 1.93 & 0.11~s & 3,266~MB & 5.04 \\
\bottomrule
\end{tabular}
\end{table}

\subsubsection{Phase 1: BCCSP Post-Quantum Signing}

Phase~1 replaces ECDSA P-256 transaction signing at the BCCSP layer with
post-quantum signatures while leaving certificates and TLS unchanged, isolating
the signature-algorithm impact. Figures~\ref{fig:phase1_write}
and~\ref{fig:phase1_mixed} present throughput-versus-send-rate curves for the
write-heavy and mixed workloads (shaded band: one-standard-deviation envelope
across three runs).

\begin{figure*}[ht]
\centering
\includesvg[inkscapelatex=false,width=\textwidth]{figures/fig_phase1_write}
\caption{Phase~1 throughput versus target send rate, write-heavy workload.}
\label{fig:phase1_write}
\end{figure*}

\begin{figure*}[ht]
\centering
\includesvg[inkscapelatex=false,width=\textwidth]{figures/fig_phase1_mixed}
\caption{Phase~1 throughput versus target send rate, mixed workload.}
\label{fig:phase1_mixed}
\end{figure*}

All four algorithms follow the same sub-knee saturation profile as ECDSA P-256,
confirming that post-quantum signing introduces no measurable overhead at
sub-saturation load. Above the knee, each post-quantum algorithm reaches a
plateau below the ECDSA ceiling, with the reduction reflecting output signature
size. Table~\ref{tab:phase1_results} presents the statistics. The apparent
write-heavy ordering FN-DSA-512 $>$ ML-DSA-65 $>$ ML-DSA-44 places ML-DSA-65
(265.6~TPS) nominally above ML-DSA-44 (263.8~TPS), but the $\sim$1.8~TPS gap is
small relative to the ML-DSA-65 cell's SD of 6.55~TPS (the largest in the study,
attributable to a single transient network event at 700~TPS in one run); this
cross-over is a within-noise inversion. Under the mixed workload the ordering
converges to the headline ranking FN-DSA-512 $>$ ML-DSA-44 $>$ ML-DSA-65,
consistent with the output-size-dominates finding. Plateau-region latency is
0.18--0.20~s for all post-quantum configurations, indistinguishable from the
0.18-s ECDSA baseline.

\begin{table}[ht]
\centering
\caption{Phase~1 throughput results (BCCSP signing only). Means across three
runs; $\Delta$~ECDSA is the signed percentage change versus the ECDSA P-256
baseline.}
\label{tab:phase1_results}
\begin{tabular}{@{}lrrrr@{}}
\toprule
\textbf{Algorithm} &
\multicolumn{2}{c}{\textbf{Write-heavy}} &
\multicolumn{2}{c}{\textbf{Mixed}} \\
\cmidrule(lr){2-3}\cmidrule(lr){4-5}
& \textbf{TPS (SD)} & \textbf{$\Delta$ ECDSA} &
  \textbf{TPS (SD)} & \textbf{$\Delta$ ECDSA} \\
\midrule
ECDSA P-256 & 275.3~(1.35) & — & 344.2~(1.93) & — \\
FN-DSA-512  & 270.7~(2.28) & $-1.7\%$ & 338.2~(0.48) & $-1.7\%$ \\
ML-DSA-44   & 263.8~(2.78) & $-4.2\%$ & 339.2~(1.50) & $-1.4\%$ \\
ML-DSA-65   & 265.6~(6.55) & $-3.5\%$ & 330.0~(2.06) & $-4.1\%$ \\
\bottomrule
\end{tabular}
\end{table}

\subsubsection{Phase 2: Post-Quantum Certificate Infrastructure}

Phase~2 adds post-quantum X.509 certificate infrastructure: the CA issues
post-quantum-signed certificates and \texttt{buildPQCChain()} validates them at
every committing peer. This isolates the certificate-processing overhead, in
particular Op-16 (certificate chain re-validation at the committing peer), which
processes the signing certificate for every endorsement signature in every
committed block. Figure~\ref{fig:phase2_throughput} presents the maximum
sustainable throughput for both workloads.

\begin{figure*}[ht]
\centering
\includesvg[inkscapelatex=false,width=\textwidth]{figures/fig_phase2_throughput}
\caption{Phase~2 maximum sustainable throughput versus ECDSA P-256 baseline, both workloads.}
\label{fig:phase2_throughput}
\end{figure*}

The certificate layer introduces substantially larger reductions than Phase~1 for
the ML-DSA family. FN-DSA-512 declines by 3.3\% write-heavy and 3.2\% mixed ---
a modest increase from the 1.7\% Phase~1 write-heavy penalty, reflecting the cost
of validating the 897-byte FN-DSA-512 SPKI in each endorsement certificate.
ML-DSA-44 declines by 9.9\% write-heavy and 8.9\% mixed; ML-DSA-65 declines by
18.1\% write-heavy and 13.0\% mixed --- the largest single-phase impact measured
(Table~\ref{tab:phase2_results}).

\begin{table}[ht]
\centering
\caption{Phase~2 throughput results (BCCSP signing plus post-quantum certificate
infrastructure). Means across three runs. Write latency is the mean transaction
latency in the write-heavy plateau region.}
\label{tab:phase2_results}
\setlength{\tabcolsep}{3.0pt}
\begin{tabular}{@{}lrrrrr@{}}
\toprule
\textbf{Algorithm} &
\multicolumn{2}{c}{\textbf{Write-heavy}} &
\multicolumn{2}{c}{\textbf{Mixed}} &
\textbf{Write lat.} \\
\cmidrule(lr){2-3}\cmidrule(lr){4-5}
& \textbf{TPS (SD)} & \textbf{$\Delta$ ECDSA} &
  \textbf{TPS (SD)} & \textbf{$\Delta$ ECDSA} & \\
\midrule
ECDSA P-256 & 275.3~(1.35) & — & 344.2~(1.93) & — & 0.18~s \\
FN-DSA-512  & 266.1~(0.35) & $-3.3\%$ & 333.2~(0.64) & $-3.2\%$ & 1.71~s \\
ML-DSA-44   & 247.9~(0.78) & $-9.9\%$ & 313.5~(2.14) & $-8.9\%$ & 5.47~s \\
ML-DSA-65   & 225.4~(0.44) & $-18.1\%$ & 299.5~(0.53) & $-13.0\%$ & 10.20~s \\
\bottomrule
\end{tabular}
\end{table}

The most significant finding is not the throughput reduction but the certificate
processing latency. Figure~\ref{fig:latency_phase} presents average transaction
latency at saturation across all phases against the 0.18-s ECDSA reference.
Phase~1 latencies are indistinguishable from baseline; the transition to Phase~2
produces a qualitative change: FN-DSA-512 rises to 1.71~s (9.5$\times$),
ML-DSA-44 to 5.47~s (30$\times$), and ML-DSA-65 to 10.20~s (57$\times$). This
reflects the certificate-validation bottleneck at Op-16: every endorsement
certificate in every block is chain-validated through the MSP trust
hierarchy~\cite{thakkar2018}, and a 500-transaction write-heavy block carries up
to 500 such validations, each requiring a full post-quantum verification over an
SPKI of 897~bytes (FN-DSA-512), 1{,}312~bytes (ML-DSA-44), or 1{,}952~bytes
(ML-DSA-65) versus 65~bytes for ECDSA. Mixed-workload latencies are far lower
(0.14, 0.18, 0.77~s) because fewer write transactions per unit time mean fewer
endorsement certificates per block.

\begin{figure*}[ht]
\centering
\includesvg[inkscapelatex=false,width=\textwidth]{figures/fig_latency_phase}
\caption{Average transaction latency at saturation by algorithm and phase, write-heavy workload.}
\label{fig:latency_phase}
\end{figure*}

\subsubsection{Phase 3: Complete \texorpdfstring{M\textsuperscript{2}Fabric}{M2Fabric} Deployment}

Phase~3 activates hybrid TLS (\texttt{X25519MLKEM768}) for all inter-node
connections, completing the deployment. Figure~\ref{fig:phase3_throughput} and
Table~\ref{tab:phase3_results} present the results. FN-DSA-512 achieves 265.5~TPS
write-heavy and 332.7~TPS mixed ($-3.5\%$, $-3.3\%$); ML-DSA-44 achieves 245.2
and 317.0~TPS ($-10.9\%$, $-7.9\%$); ML-DSA-65 achieves 225.8 and 301.1~TPS
($-18.0\%$, $-12.5\%$).

\begin{figure*}[ht]
\centering
\includesvg[inkscapelatex=false,width=\textwidth]{figures/fig_phase3_throughput}
\caption{Phase~3 complete \mfabric{} deployment throughput versus ECDSA P-256 baseline, both workloads.}
\label{fig:phase3_throughput}
\end{figure*}

\begin{table}[ht]
\centering
\caption{Phase~3 throughput results (complete \mfabric{} deployment: BCCSP,
certificate infrastructure, and hybrid TLS). Means across three runs.}
\label{tab:phase3_results}
\setlength{\tabcolsep}{3.5pt}
\begin{tabular}{@{}lrrrrr@{}}
\toprule
\textbf{Algorithm} &
\multicolumn{2}{c}{\textbf{Write-heavy}} &
\multicolumn{2}{c}{\textbf{Mixed}} &
\textbf{Write lat.} \\
\cmidrule(lr){2-3}\cmidrule(lr){4-5}
& \textbf{TPS (SD)} & \textbf{$\Delta$ ECDSA} &
  \textbf{TPS (SD)} & \textbf{$\Delta$ ECDSA} & \\
\midrule
ECDSA P-256 & 275.3~(1.35) & — & 344.2~(1.93) & — & 0.18~s \\
FN-DSA-512  & 265.5~(1.42) & $-3.5\%$ & 332.7~(4.12) & $-3.3\%$ & 1.80~s \\
ML-DSA-44   & 245.2~(0.73) & $-10.9\%$ & 317.0~(2.14) & $-7.9\%$ & 5.49~s \\
ML-DSA-65   & 225.8~(2.09) & $-18.0\%$ & 301.1~(2.88) & $-12.5\%$ & 10.52~s \\
\bottomrule
\end{tabular}
\end{table}

Across the six cells the Phase~2-to-Phase~3 change is small and spans zero, and
every difference lies within the corresponding Phase~3 standard-deviation
envelope. The hybrid-TLS layer therefore contributes a per-transaction overhead
below the detection threshold of the $n=3$ design: connection-level key-exchange
cost is amortised across many transactions per TLS session and adds no
significant load on a network already bottlenecked by certificate validation.

\subsection{Consolidated Analysis and Performance Models}
\label{sec:res_consolidated}

Figure~\ref{fig:phase_progression} presents the throughput trajectory for each
algorithm across all three phases. The progression reveals a structurally
different cost profile across families (attribution in
Table~\ref{tab:phase_attribution}): for FN-DSA-512 the 3.5\% write-heavy
reduction splits roughly evenly between Phases~1 and~2; for ML-DSA-44 and
especially ML-DSA-65 the certificate layer dominates, adding 5.7 and 14.6
percentage points respectively, while the Phase~3 TLS layer contributes under 0.2
points in every case. This confirms that the dominant determinant is not signing
speed but the per-endorsement payload --- 1{,}563~bytes for FN-DSA-512, 3{,}732
for ML-DSA-44, and 5{,}261 for ML-DSA-65 --- applied to every endorsement in
every block, independently corroborating the earlier PQFabric finding of Holcomb
et al.~\cite{holcomb2021}.

\begin{figure*}[ht]
\centering
\includesvg[inkscapelatex=false,width=\textwidth]{figures/fig_phase_progression}
\caption{Phase-by-phase throughput progression across the three \mfabric{} deployment stages.}
\label{fig:phase_progression}
\end{figure*}

\begin{table*}[ht]
\centering
\caption{Phase-by-phase throughput attribution. Values show absolute TPS and the
contribution of each phase to the total reduction from the ECDSA P-256 baseline.
Phase~3 contribution is the change from Phase~2 to Phase~3 (hybrid TLS overhead).}
\label{tab:phase_attribution}
\setlength{\tabcolsep}{3.5pt}
\begin{tabular}{@{}lcrrrrr@{}}
\toprule
\textbf{Algorithm} & \textbf{Workload} &
\textbf{ECDSA} & \textbf{Phase 1} & \textbf{Phase 2} &
\textbf{Phase 3} & \textbf{Total $\Delta$} \\
\midrule
\multirow{2}{*}{FN-DSA-512}
  & Write & 275.3 & 270.7~($-1.7\%$) & 266.1~($-1.6$pp) & 265.5~($-0.2$pp) & $-3.5\%$ \\
  & Mixed & 344.2 & 338.2~($-1.7\%$) & 333.2~($-1.5$pp) & 332.7~($-0.2$pp) & $-3.3\%$ \\
\addlinespace
\multirow{2}{*}{ML-DSA-44}
  & Write & 275.3 & 263.8~($-4.2\%$) & 247.9~($-5.7$pp) & 245.2~($-1.0$pp) & $-10.9\%$ \\
  & Mixed & 344.2 & 339.2~($-1.4\%$) & 313.5~($-7.5$pp) & 317.0~($+1.1$pp) & $-7.9\%$ \\
\addlinespace
\multirow{2}{*}{ML-DSA-65}
  & Write & 275.3 & 265.6~($-3.5\%$) & 225.4~($-14.6$pp) & 225.8~($+0.2$pp) & $-18.0\%$ \\
  & Mixed & 344.2 & 330.0~($-4.1\%$) & 299.5~($-8.9$pp) & 301.1~($+0.5$pp) & $-12.5\%$ \\
\bottomrule
\multicolumn{7}{@{}l@{}}{\footnotesize pp = percentage points, relative to prior phase.} \\
\end{tabular}
\end{table*}

\subsubsection{Multiplicative Throughput Degradation Model}

The attribution data admit a multiplicative decomposition. Let
$\delta_p(\mathrm{alg})$ denote the fractional throughput reduction introduced by
phase $p\in\{1,2,3\}$ relative to the immediately preceding phase. The Phase~3
throughput and total degradation are

\begin{equation}
\Gamma_{\mathrm{Ph3}}
= \Gamma_{\mathrm{ECDSA}}
  \bigl(1-\delta_1\bigr)\bigl(1-\delta_2\bigr)\bigl(1-\delta_3\bigr),
\label{eq:tps_degradation}
\end{equation}

\begin{equation}
\delta_{\mathrm{total}}
= 1-\bigl(1-\delta_1\bigr)\bigl(1-\delta_2\bigr)\bigl(1-\delta_3\bigr)
\;\approx\; \delta_1+\delta_2+\delta_3
\quad\text{for small }\delta_p,
\label{eq:phase_attr}
\end{equation}

where $\Gamma_{\mathrm{ECDSA}}$ is the baseline throughput, the $(1-\delta_p)$
factors are the per-phase retention ratios, and the additive approximation on the
right of~(\ref{eq:phase_attr}) holds when each per-phase reduction $\delta_p$ is
small. Table~\ref{tab:degradation_model} applies these to the write-heavy values.
For FN-DSA-512 ($\max\delta_p=0.017$) the approximation is exact to three
decimals; for ML-DSA-65 ($\delta_2=0.146$) the additive-vs-exact gap widens to
0.5~pp. The approximation should be used only when all per-phase reductions are
below $\sim$0.05; otherwise the exact multiplicative form is preferred.

\begin{table}[ht]
\centering
\caption{Multiplicative degradation model applied to write-heavy Phase~3 results.
$\delta_1,\delta_2,\delta_3$ are per-phase fractional reductions (negative
denotes a small phase-to-phase gain from run-to-run variation).
$\delta_{\mathrm{total}}^{\mathrm{mult}}$ is the exact result;
$\delta_{\mathrm{total}}^{\mathrm{add}}$ the additive approximation.}
\label{tab:degradation_model}
\setlength{\tabcolsep}{4pt}
\footnotesize
\begin{tabular}{@{}lrrrrrrr@{}}
\toprule
\textbf{Algorithm} &
\textbf{$\delta_1$} & \textbf{$\delta_2$} & \textbf{$\delta_3$} &
\textbf{$\delta_{\mathrm{total}}^{\mathrm{mult}}$} &
\textbf{$\delta_{\mathrm{total}}^{\mathrm{add}}$} &
\textbf{Measured} & \textbf{Model err.} \\
\midrule
FN-DSA-512 & 0.017 & 0.016 & 0.002 & 0.035 & 0.035 & $-3.5\%$ & $<$0.1\,pp \\
ML-DSA-44  & 0.042 & 0.057 & 0.010 & 0.106 & 0.109 & $-10.9\%$ & 0.3\,pp \\
ML-DSA-65  & 0.035 & 0.146 & $-0.002$ & 0.178 & 0.179 & $-18.0\%$ & 0.2\,pp \\
\bottomrule
\end{tabular}
\end{table}

\subsubsection{Statistical Confidence Intervals}

All throughput figures are means across three independent runs ($n=3$). The 95\%
confidence interval for the true mean $\mu_\Gamma$ given the sample mean
$\hat{\Gamma}$ and sample standard deviation $s$ is

\begin{equation}
\hat{\Gamma}\;\pm\;t_{0.025,\,2}\,\frac{s}{\sqrt{3}},
\label{eq:tps_ci}
\end{equation}

where $t_{0.025,2}=4.303$ is the upper 2.5th percentile of the Student-$t$
distribution with $n-1=2$ degrees of freedom and $s/\sqrt{3}$ is the standard
error of the mean. The intervals of adjacent algorithms are non-overlapping at
95\% for the headline write-heavy comparisons, supporting the ordering
FN-DSA-512 $>$ ML-DSA-44 $>$ ML-DSA-65. The $n=3$ sample size is not designed to
detect small effects, but the headline ordering is robust on three independent
checks: pairwise non-overlapping Student-$t$ intervals; a non-parametric bootstrap
($10^4$ resamples) giving tighter but order-preserving Phase~3 write-heavy
intervals; and very large inter-algorithm effect sizes (Cohen's $d\approx18$ and
$\approx12$ against the conventional ``large'' threshold of 0.8). Only the
small-effect comparisons (notably the hybrid-TLS contribution) are underpowered;
future work should re-run at $n\ge5$ with analysis-of-variance.

\subsubsection{Certificate Validation Latency Model}

The Phase~2 write latencies are explained by the certificate-validation
bottleneck at Op-16. Let $T_b$ be the number of transactions per block,
$\bar{e}$ the mean endorsers per transaction, and
$\tau_{\mathrm{cert}}(\mathrm{alg})$ the per-certificate validation time. The
additional per-block validation time introduced by Phase~2 relative to Phase~1 is

\begin{equation}
\Delta t_{\mathrm{block}}(\mathrm{alg})
= T_b\,\bar{e}\,
\bigl[\tau_{\mathrm{cert}}(\mathrm{alg})-\tau_{\mathrm{cert}}(\mathrm{ECDSA})\bigr],
\label{eq:cert_latency}
\end{equation}

where $T_b$, $\bar{e}$, and $\tau_{\mathrm{cert}}$ are as defined above. The
average transaction latency at saturation is proportional to
$\Delta t_{\mathrm{block}}$ when certificate validation is the binding
constraint. The empirical latencies are consistent with $\tau_{\mathrm{cert}}$
scaling super-linearly with SPKI size: the ML-DSA-65-to-FN-DSA-512 latency ratio
$10.20/1.71\approx6.0$ exceeds the public-key size ratio
$1{,}952/897\approx2.18$. A power-law fit
$\tau_{\mathrm{cert}}\propto|\mathit{pk}|^{\beta}$ gives
$\beta\approx2.30$, indicating that doubling the public-key size approximately
quintuples the per-certificate validation time.

\subsubsection{Ledger Storage Overhead Model and Secondary Metrics}

The near-unity storage multipliers in Table~\ref{tab:storage_results} follow from
a simple block-payload model. Let $N_{\mathrm{tx}}$ be the total committed
transactions, $\bar{e}$ the mean endorsers per transaction, and
$|\sigma_{\mathrm{alg}}|$ the per-endorsement signature size. The incremental
signature storage relative to ECDSA and the resulting storage multiplier are

\begin{equation}
\Delta S(\mathrm{alg})
= N_{\mathrm{tx}}\,\bar{e}\,
\bigl(|\sigma_{\mathrm{alg}}|-|\sigma_{\mathrm{ECDSA}}|\bigr),
\label{eq:storage_model}
\end{equation}

\begin{equation}
\mu_S(\mathrm{alg})
= \frac{S_{\mathrm{ECDSA}}+\Delta S(\mathrm{alg})}{S_{\mathrm{ECDSA}}},
\label{eq:storage_multiplier}
\end{equation}

where $S_{\mathrm{ECDSA}}$ is the baseline ledger size and the remaining symbols
are as above. For the write-heavy benchmark
($N_{\mathrm{tx}}\approx10^5$, $\bar{e}=1$, $S_{\mathrm{ECDSA}}=1{,}690$~MB),
$\Delta S(\text{FN-DSA-512})\approx60$~MB gives $\mu_S\approx1.04$ from the model,
consistent with the measured $0.99\times$. The model confirms that the dominant
storage component remains the transaction payload, not the signature field.

Figure~\ref{fig:storage} and Table~\ref{tab:storage_results} present the ledger
storage per peer. All three post-quantum algorithms produce write-heavy ledgers
within 4\% of the 1{,}690~MB ECDSA baseline (0.97--0.99$\times$). The
mixed-workload figures need care: the network was not reset between the three
sequential runs, so each snapshot is cumulative; normalising the run~1 snapshot
against the ECDSA run~1 figure gives $\sim$0.86$\times$, reflecting the lower
block count at reduced throughput rather than any per-transaction saving.

\begin{figure*}[ht]
\centering
\includesvg[inkscapelatex=false,width=\textwidth]{figures/fig_storage}
\caption{Ledger storage per peer by algorithm and workload, measured after each benchmark run.}
\label{fig:storage}
\end{figure*}

\begin{table}[ht]
\centering
\caption{Ledger storage results. Write-heavy values are single-run snapshots from
the Phase~1 benchmark. Mixed values are first-run snapshots normalised against the
ECDSA P-256 first-run figure of 3{,}565~MB.}
\label{tab:storage_results}
\begin{tabular}{@{}lrrrr@{}}
\toprule
\textbf{Algorithm} &
\multicolumn{2}{c}{\textbf{Write-heavy (single run)}} &
\multicolumn{2}{c}{\textbf{Mixed (per-run estimate)}} \\
\cmidrule(lr){2-3}\cmidrule(lr){4-5}
& \textbf{MB/peer} & \textbf{Multiplier} &
  \textbf{MB/peer (run~1)} & \textbf{Multiplier} \\
\midrule
ECDSA P-256 & 1,690 & 1.00$\times$ & 3,565 & 1.00$\times$ \\
FN-DSA-512  & 1,669 & 0.99$\times$ & 3,070 & 0.86$\times$ \\
ML-DSA-44   & 1,633 & 0.97$\times$ & 2,822 & 0.79$\times$ \\
ML-DSA-65   & 1,679 & 0.99$\times$ & 3,005 & 0.84$\times$ \\
\bottomrule
\end{tabular}
\end{table}

Table~\ref{tab:block_prop} presents block propagation metrics from Prometheus.
All post-quantum configurations achieve 9.5--9.8~blocks/s mixed and
17.4--18.4~blocks/s write-heavy. Average block processing time is uniformly
$\sim$1.0--1.1~ms for all configurations except the ML-DSA-44 mixed measurement
(high variance from a warm-up transient in one run). The $\sim$1~ms per-block
figure confirms that the Phase~2/3 reductions are driven by certificate-validation
overhead at the committing peer rather than gossip propagation or orderer latency.

\begin{table}[ht]
\centering
\caption{Block propagation metrics from sustained-load Prometheus measurements.
Means across runs where multiple measurements are available.}
\label{tab:block_prop}
\begin{tabular}{@{}lrrr@{}}
\toprule
\textbf{Configuration} & \textbf{Workload} &
\textbf{Blocks/s} & \textbf{Avg proc.\ time} \\
\midrule
ECDSA P-256 & Mixed  & 5.04 & 1.13~ms \\
\addlinespace
FN-DSA-512  & Write  & 18.19 & 1.02~ms \\
FN-DSA-512  & Mixed  & 9.69  & 1.06~ms \\
ML-DSA-44   & Write  & 17.94 & 1.27~ms \\
ML-DSA-44   & Mixed  & 9.74  & 1.04~ms \\
ML-DSA-65   & Write  & 18.03 & 1.02~ms \\
ML-DSA-65   & Mixed  & 9.51  & 1.07~ms \\
\bottomrule
\end{tabular}
\end{table}

\subsection{Algorithm Selection Guidance}
\label{sec:res_selection}

The phase attribution confirms the prediction of the cost model
(eq.~\ref{eq:selection_cost}): per-endorsement payload size, not signing speed,
is the dominant determinant. FN-DSA-512 delivers the highest end-to-end
throughput despite signing $\sim$53\% more slowly than ML-DSA-44 in
microbenchmarks, because its smaller signature and 897-byte public key give a far
smaller per-endorsement payload. Table~\ref{tab:selection_matrix} provides the
selection matrix. FN-DSA-512 is recommended where throughput headroom is
constrained or for early-transition deployments; ML-DSA-44 for Category~2
assurance where a $\sim$10\% reduction is acceptable; and ML-DSA-65 for
highest-assurance environments mandating Category~3, with the 18\% reduction and
tenfold latency increase explicitly acknowledged in the deployment design.

\begin{table*}[ht]
\centering
\caption{Algorithm selection matrix. Throughput and latency figures are
Phase~3 results from Tables~\ref{tab:phase3_results}
and~\ref{fig:latency_phase}. Security categories follow NIST FIPS~204
(finalised 2024) and the FIPS~206 Initial Public Draft (2025). $\Delta_{\mathrm{ECDSA}}$
is the throughput change relative to the ECDSA P-256 baselines of 275.3~TPS
write-heavy and 344.2~TPS mixed. The latency multiplier $L/L_{\mathrm{ECDSA}}$
divides the Phase~3 write-heavy latency by the 0.18~s ECDSA reference.}
\label{tab:selection_matrix}
\setlength{\tabcolsep}{3pt}
\footnotesize
\begin{tabular}{@{}lcrrrrcp{3.4cm}@{}}
\toprule
\textbf{Algorithm} & \textbf{Sec.\ Cat.} &
\multicolumn{2}{c}{\textbf{Write-heavy}} &
\multicolumn{2}{c}{\textbf{Mixed}} &
\textbf{$L/L_{\mathrm{ECDSA}}$} &
\textbf{Recommendation} \\
\cmidrule(lr){3-4}\cmidrule(lr){5-6}
& & \textbf{TPS} & \textbf{$\Delta_{\mathrm{ECDSA}}$} &
\textbf{TPS} & \textbf{$\Delta_{\mathrm{ECDSA}}$} & & \\
\midrule
FN-DSA-512 & 1 & 265.5 & $-3.5\%$  & 332.7 & $-3.3\%$  & $10\times$  & Primary: lowest overhead; recommended for throughput-sensitive deployments and early-transition pilots \\
ML-DSA-44  & 2 & 245.2 & $-10.9\%$ & 317.0 & $-7.9\%$  & $30\times$  & Balanced: NIST first-choice signature; recommended where Cat.~2 assurance is mandated and $\sim$11\% reduction is acceptable \\
ML-DSA-65  & 3 & 225.8 & $-18.0\%$ & 301.1 & $-12.5\%$ & $58\times$  & High-assurance: recommended only where Cat.~3 is mandated; 18\% reduction and tenfold latency increase must be acknowledged in the deployment design \\
\bottomrule
\end{tabular}
\end{table*}

\subsection{Comparison with Prior Integrations}
\label{sec:res_comparison}

Revisiting the comparison of Section~\ref{sec:litreview} in light of the results,
\mfabric{} addresses all four live cryptographic layers using only standard Go
tooling and FIPS-standardised algorithms, extends the live CA for runtime-enrolled
identities, and provides the \texttt{buildPQCChain()} solution to the Go X.509
barrier. No prior work addresses all of these simultaneously, and the components
\mfabric{} addresses constitute the full set required by the formal integration
requirements (eq.~\ref{eq:coverage_constraint}). Beyond scope, the work
contributes the first end-to-end performance characterisation of a four-layer
FIPS-standardised integration on the current LTS; the finding that
per-endorsement payload size dominates throughput independently corroborates and
extends the earlier PQFabric observation to the FIPS-finalised generation. The
choice of pure-Go over CGo bindings~\cite{stebila2016oqs} also distinguishes the
work: pure Go needs no C toolchain, avoids cross-language memory management, and
keeps the constant-time properties off the CGo boundary, while CIRCL's
production deployment and independent review give a reliability baseline
comparable to \texttt{liboqs}.

\subsection{Security Analysis of the Implementation}
\label{sec:res_security}

The security rests on three foundations: the security of the post-quantum
algorithms (established by NIST's six-year standardisation
process~\cite{alagic2022nist}; no polynomial-time quantum algorithm is known for
the underlying Module LWE or NTRU/SIS problems), the correctness of the
implementation (established by the verification of Section~\ref{sec:exp_verify}),
and the integrity of the integration design (which preserves the existing Fabric
security model: the BCCSP contract is satisfied exactly, MSP policy evaluation is
unchanged, the TLS configuration is a subset of Go's option space, and CA-issued
certificates follow the RFC~5280 structure).

The custom code paths added above the primitive layer are not covered by upstream
review and warrant explicit identification of residual attack surface. Four
classes of concern apply. \textit{Malformed-certificate handling}: the custom
ASN.1 parser in \texttt{buildPQCChain()} operates on attacker-controlled input;
it rejects parsing failures by error-return rather than panic, but
positive evidence of robustness requires coverage-guided fuzz testing, which
is not yet conducted. \textit{Algorithm-OID spoofing}: the dispatch table checks
issuer-SPKI/child consistency, so a substituted classical-for-post-quantum
signature fails; however, certificates presenting a known PQC OID with
wrong-length key bytes have not been tested under adversarial corpora.
\textit{Downgrade resistance}: the legacy SHA-256-prehash fallback introduces a
second acceptance path; a strict-mode toggle disabling it after a cutover date is
recommended for production. \textit{Hybrid-TLS partial coverage}:
\texttt{X25519MLKEM768} provides session-key confidentiality but not
quantum-resistant TLS-certificate authentication, which remains classical under
the dual-CA design, so TS4 is closed only on the confidentiality dimension.

\subsubsection{Planned Adversarial-Evaluation Workstream}
\label{sec:res_adversarial_plan}

The fuzz-testing and adversarial-corpus gaps identified above are addressed
by a four-stage adversarial evaluation scheduled as the immediate follow-up
to the present work. Stage~1 runs a \texttt{go-fuzz}-based coverage-guided
campaign against the \texttt{buildPQCChain()} ASN.1 entry points, with a
seed corpus drawn from valid PQC chains in the test network and a target
of $\geq 24$~CPU-hours per algorithm under address-sanitizer
instrumentation; success is defined as no panic, no out-of-bounds read, and
no signature-validation false positive across the full campaign.
Stage~2 constructs an adversarial corpus exercising the OID-spoofing
surface: certificates with valid PQC OIDs and truncated, oversized, or
zero-byte SPKI payloads; cross-algorithm OID/key mismatches; AKID/SKID
disjoint chains; and validity-window edge cases at the
\texttt{notBefore}/\texttt{notAfter} boundary. Stage~3 evaluates downgrade
resistance by toggling the legacy SHA-256-prehash fallback off and
confirming that all post-fix-cert-generator chains continue to verify
while all pre-fix ones are rejected as intended after the cutover.
Stage~4 establishes side-channel observability against the deployed
binaries using a constant-time differential analysis on the signing and
verification paths, looking for input-dependent timing variations that
would indicate a regression in the upstream constant-time guarantees of
CIRCL or \texttt{go-fn-dsa}. The output of the workstream is a public
adversarial-evaluation report releasable alongside the source tree, which
moves the implementation from \emph{functionally verified} to
\emph{adversarially robust}.

These surfaces do not invalidate the positive functional verification but
establish that the implementation's security stance is \textit{provisional}
pending the adversarial evaluation described in
Section~\ref{sec:res_adversarial_plan} and the limitations consolidated
in Section~\ref{sec:disc_limits}.
The residual TLS-CA quantum risk is bounded, time-limited, and active-attack-gated
(a man-in-the-middle on future sessions requiring an active network position),
materially different from the passive HNDL-style TS1 attack on
permanently-recorded ledger material, and is properly characterised as a
transitional limitation rather than a fundamental compromise.

%=====================================================================
\section{Discussion} \label{sec:discussion}
%=====================================================================

This section consolidates the implications of the results of
Section~\ref{sec:experimental} along four dimensions: the completeness of
the multilayer approach relative to prior partial integrations, the
amplification of the harvest-now-decrypt-later threat by ledger
immutability, the generalisability of the architecture to other permissioned
blockchain platforms, and the alignment of the design with emerging national
and international post-quantum migration strategies.

\subsection{Completeness of the Multilayer Approach}
\label{sec:disc_completeness}

The integration-coverage identity~(\ref{eq:integration_coverage}) gives
$\delta_{\mfabric{}}=4/4=1$ with respect to the four ECDLP-vulnerable layers
identified by the threat analysis. The significance of this is structural
rather than cosmetic. Equation~(\ref{eq:residual_surface}) shows that the
residual quantum-attack surface $R(\phi)$ is a weighted sum over the
threats $\mathrm{TS}_i$ closed by phase $\phi$, with the weight $s_i$
determined by severity. Any integration that leaves a strict subset of the
four layers classical maintains a strictly positive residual surface; in
particular, a BCCSP-only integration (the implicit shape of several prior
proposals) leaves TS1 entirely open, because a quantum adversary who
recovers the classical CA's private key can still issue fraudulent
\emph{post-quantum} certificates that the post-quantum MSP accepts as
legitimate. The multilayer design exists precisely to drive $R(\phi)$ to
zero, and the empirical Phase-by-Phase progression in
Table~\ref{tab:phase_attribution} confirms that each successive phase
contributes a real, measurable, and bounded throughput cost---i.e.\ that
the coverage gain is not paid for by performance collapse.

The most operationally consequential implication is that
\emph{partial migration is a security anti-pattern, not a transitional
state}: an operator who deploys only Phase~1 obtains no protection against
the highest-severity threat (TS1), while paying most of the throughput
cost. This finding extends the architectural recommendation of the
underlying threat analysis from a normative claim to an empirically
defensible one.

\subsection{Ledger Immutability and HNDL Amplification}
\label{sec:disc_hndl}

A central conceptual finding is that permissioned blockchains amplify the
HNDL threat in a structurally different way from conventional PKI. Unlike a
classical TLS deployment, where key rotation and migration eventually
reduce the harvestable set, the append-only ledger
guarantees that $|\mathcal{H}(t)|$ is non-decreasing, so every
ECDSA-signed transaction, endorsement, block header, and certificate
committed before $t_m$ remains a permanent harvest target even after
migration completes. Equation~(\ref{eq:integration_coverage}) addresses
\emph{future} cryptographic material; it cannot retroactively remove
historical ECDLP-bound material from the chain.

The practical implication is that organisations operating long-lived Fabric
networks should treat \mfabric{} not as protection against an abstract future
adversary but as a mechanism that bounds the size of the permanent harvest
set. Each day of continued ECDSA operation enlarges $|\mathcal{H}|$ by
$\bar{\Gamma}\bar{k}\cdot 86{,}400$ keys, where $\bar{\Gamma}$ is mean
throughput and $\bar{k}$ is the mean number of new ECDSA keys per
transaction. For the test network's 275.3~TPS write-heavy plateau this
amounts to $\sim 2.4\times 10^{7}$ key-equivalents per day. Migration
urgency is therefore proportional not only to the residual security
lifetime $Z$ in Mosca's inequality but to the rate at which the
unprotected ledger continues to accumulate harvestable material.

\subsection{Generalisability to Other Permissioned Blockchains}
\label{sec:disc_general}

The four-layer decomposition---signing, certificate infrastructure,
identity validation, transport---is not Fabric-specific. The same
operational categories recur in Quorum, Besu, Corda, and most
enterprise-grade permissioned platforms: a per-node signing module, a
certificate authority, a participant-validation surface, and a TLS
transport. The cryptographic-agility, dispatch-table, and OID-aware
chain-validation patterns developed for \mfabric{} therefore transfer
directly. The most transferable single artefact is the
\texttt{buildPQCChain()} pattern itself: any Go-based blockchain platform
that imports \texttt{crypto/x509} is currently subject to the same
post-quantum OID rejection, and the same dispatch-on-\texttt{signatureAlgorithm}
structure resolves it.

Two qualifications apply. First, FN-DSA-512's OID is provisionally assigned
within the OQS experimental arc \texttt{1.3.9999.3.x}; until FIPS~206 is
finalised and a stable OID is published, any portable validator must support
the same transition-window dispatch \mfabric{} implements. Second, hybrid
TLS support is currently a Go-standard-library feature
(\texttt{tls.X25519MLKEM768}, Go~1.24+); platforms targeting other language
ecosystems must either wait for analogous stdlib support or accept the
maintenance cost of in-process key-exchange code.

\subsection{Alignment with National and International PQC Migration Strategies}
\label{sec:disc_policy}

The \mfabric{} design choices align closely with three concurrently
developing migration frameworks: the NIST transition
roadmap~\cite{nist_ir8547}, which recommends RSA/ECC deprecation by 2030
and disallowance by 2035; the NSA Commercial National Security Algorithm
Suite~2.0~\cite{nsa_cnsa2}, with full adoption mandated by 2035; and
emerging national frameworks such as the draft \emph{Brunei
Post-Quantum Cryptography Migration Framework}~\cite{csb_brunei2026}, which
adopts the same 2030/2035 anchors and an explicit five-phase
Assess--Select--Validate--Deploy--Monitor strategy.

The mapping to the latter framework is direct: \mfabric{}'s three-phase
migration (BCCSP~$\to$~CA+MSP~$\to$~hybrid TLS) sequences the framework's
Validate and Deploy phases for the Fabric platform specifically; the
algorithm selections coincide with the framework's recommended set
(FIPS~203/204/205 final, FIPS~206 draft); and the empirical Phase-3
characterisation in Section~\ref{sec:res_consolidated} supplies the
quantitative pilot evidence the framework's Phase~3 (Validation and Pilot)
explicitly requires. \mfabric{} is therefore complementary to these
policy-level frameworks rather than competitive with them: the frameworks
specify the \emph{what} and \emph{when} at national or international
policy level; \mfabric{} supplies a concrete \emph{how} for the
permissioned-blockchain layer of the migration.

\subsection{Limitations}
\label{sec:disc_limits}

Five limitations bound the present results.

First, the dual-CA architecture is a transitional design: the residual
TLS certificate authentication path remains classical until
\texttt{crypto/tls} accepts post-quantum certificates, leaving an
active-attack-gated residual exposure that is materially weaker than a
passive HNDL exposure but not zero.

Second, the empirical evaluation uses a single-host, two-organisation
network ($n=3$) selected to isolate cryptographic overhead from network
latency; throughput on multi-host deployments will be additionally bounded
by gossip-propagation costs that scale with signature size, and a
multi-host re-run with $n\geq 5$ is required to characterise this regime.

Third, verification establishes functional rather than adversarial
robustness: coverage-guided fuzz testing of the custom ASN.1 parser, an
adversarial corpus targeting OID-spoofing edge cases, and a strict-mode
toggle disabling the legacy SHA-256-prehash fallback are scheduled for the
follow-up adversarial-evaluation workstream.

Fourth, the FN-DSA-512 OID assignment is provisional; the transition to
the finalised FIPS~206 OID is a single-constant change at the source level
but requires a transition-window entry in the validator's dispatch table to
preserve verifiability of certificates issued under the experimental
assignment.

Fifth, side-channel resistance at the primitive level depends on the
upstream constant-time guarantees of CIRCL and \texttt{go-fn-dsa};
deployment in HSM-backed environments and characterisation against
electromagnetic and fault-injection adversaries~\cite{ravi2024sidechannel}
remain open and are properly the subject of dedicated follow-up work.

\subsection{Threats to Validity}
\label{sec:disc_threats}

Three classes of validity threat warrant explicit identification.
\emph{Internal validity}: the $n=3$ sample size is underpowered for
small-effect detection (notably the hybrid-TLS Phase~3 contribution); the
headline ordering is supported by non-overlapping Student-$t$ intervals,
a non-parametric bootstrap, and very large Cohen's-$d$ values, but smaller
effects should be re-examined at $n\geq 5$ with analysis-of-variance.
\emph{External validity}: results are reported for a single hardware
configuration and a single chaincode workload (Fabric's standard
asset-transfer-basic); generalising to multi-host, multi-chaincode, and
geographically distributed deployments requires the multi-host extension.
\emph{Construct validity}: the maximum-sustainable-throughput metric is
operationalised via the plateau-region mean in~(\ref{eq:plateau}); the
choice of saturation threshold $\varepsilon$ admits a small degree of
analyst discretion, mitigated here by reporting the per-rate achieved
throughput alongside the plateau aggregate so the reader can reconstruct
$\Gamma_{\max}$ under alternative thresholds.

%=====================================================================
\section{Conclusion} \label{sec:conclusion}
%=====================================================================
This article presented \mfabric{}, a multilayer post-quantum architecture for
Hyperledger Fabric that addresses all four live cryptographic layers
simultaneously: the BCCSP signing layer (ML-DSA-44, ML-DSA-65, FN-DSA-512), the
certificate infrastructure (a post-quantum Fabric CA with custom ASN.1
construction), the MSP layer (the \texttt{buildPQCChain()} validator), and the TLS
layer (\texttt{X25519MLKEM768} hybrid key exchange), deployed through a three-phase
migration. The principal enabler, \texttt{buildPQCChain()}, resolves the Go
standard-library barrier that has prevented live CA integration in prior work; to
the authors' knowledge this is the first runtime integration of the four primary
layers of Fabric v2.5.x using FIPS-standardised algorithms in standard Go.
Empirically, the complete deployment reduces maximum sustainable write-heavy
throughput by 3.5\% for FN-DSA-512, 10.9\% for ML-DSA-44, and 18.0\% for ML-DSA-65
relative to ECDSA P-256. The central finding is that per-endorsement payload size,
not signing speed, dominates throughput: FN-DSA-512 leads despite slower signing,
making certificate-and-signature size a first-class selection criterion. The
certificate layer (Phase~2) is the dominant cost driver, while hybrid TLS
(Phase~3) is effectively free.

These results are bounded by the transitional dual-CA design, a single-host
two-organisation benchmark ($n=3$), and verification establishing functional
rather than adversarial robustness. Future work should pursue a multi-host
benchmark at $n\ge5$ and a four-stage adversarial workstream, moving \mfabric{}
toward a deployment-ready, adversarially-evaluated distribution.

\bibliographystyle{ieeetr}
\bibliography{Ref}

\vfill\pagebreak

\end{document}